%% 3-10 ボールはいつその高さに届く？：2次方程式（第12回編集委員会：組 V3-5）
%% 流れ：はてな（高さ15 mの窓）→ やってみよう（表で高さをさがす）→ ① 2次方程式 → ② x^2=k → ③ かけて0
%%   → ④ 因数分解（くくる・和と積・x^2-a^2・くくってから）→ ⑤ 因数分解で解く（はてなの答え）→ ⑥ 平方完成で解く
%%   → ⑦ √ の計算 → ⑧ 解の公式 → ⑨ 解とグラフ・判別式 → ⑩ 式が＋になるのはどこ？
%% 第5巻の思い出そう（5-5・5-7・5-11・5-13）と同じ例を、例題・じぶんでとこうに入れた。
%% 3-8（展開・くくる・平方完成）は別の組の章。使うときは「（3-8）」と1行で思い出させる。
\makeatletter
% 式を、まん中に1行で置く  \bkthrkline{…}
\newcommand\bkthrkline[1]{\par\addvspace{3pt}{\centering #1\par}\addvspace{3pt}}
% 説明の段の見出し（丸数字＋太字の1行）  \bkthrkstep{番号}{見出し}
\newcommand\bkthrkstep[2]{\par\addvspace{6pt}\Needspace{4\baselineskip}\noindent
  {\bkcircled[bkOrange]{#1}\enspace\sffamily\bfseries\color{bkOrange!75!black}#2}\par\nobreak\vspace{2pt}}
% 本文の中の小見出し（それだけで1行）  \bkthrkhead{見出し}
\newcommand\bkthrkhead[1]{\par\addvspace{4pt}\Needspace{5\baselineskip}\noindent\textbf{#1}\par\nobreak}
% ボール  \bkthrkball{(x,y)}{半径}
\newcommand\bkthrkball[2]{\shade[ball color=bkOrange] #1 circle (#2);}
% 面積図のタイル  \bkthrktile{(左下)}{はば}{高さ}{色}{文字}
\newcommand\bkthrktile[5]{\filldraw[fill=#4!22, draw=#4!75!black, line width=0.6pt] #1 rectangle ++(#2,#3);
  \node[font=\footnotesize, text=bkInk] at ($#1+(#2/2,#3/2)$) {#5};}
% じぶんでとこうを2つ横にならべる  \bkthrkpair{グループ名}{左}{右}
\newcommand\bkthrkpair[3]{\begingroup\tcbset{equal height group=bkthrk#1}%
  \noindent\begin{minipage}[t]{0.49\linewidth}\raggedright#2\end{minipage}\hfill
  \begin{minipage}[t]{0.49\linewidth}\raggedright#3\end{minipage}\par\endgroup}
\makeatother

\bkchapter{3}{10}{ボールはいつその\R{高}{たか}さに\\\R{届}{とど}く？：2\R{次方程式}{じ|ほう|てい|しき}}{\bklevelrange{2}{3}}{ボールはいつその\R{高}{たか}さに\R{届}{とど}く？：2\R{次方程式}{じ|ほう|てい|しき}}
\bkmichishirube{2-8,2-16,3-8,3-9}{5-5,5-7,5-11,6-12}

\begin{goal}
\begin{itemize}[label={\color{\bkTT}●}, leftmargin=1.6\zw, itemsep=0pt, topsep=1pt]
\item \bklevel{2}\enspace $x^2=k$ の\R{形}{かたち}や、\R{因数分解}{いん|すう|ぶん|かい}（くくる・\R{和}{わ}と\R{積}{せき}・$x^2-a^2$）を\R{使}{つか}って、2\R{次方程式}{じ|ほう|てい|しき}を\R{解}{と}ける。
\item \bklevel{3}\enspace \R{平方完成}{へい|ほう|かん|せい}から\R{解}{かい}の\R{公式}{こう|しき}を\R{作}{つく}り、どんな2\R{次方程式}{じ|ほう|てい|しき}でも\R{解}{と}ける。$\sqrt{\phantom{2}}$ の\R{計算}{けい|さん}ができる。
\item \bklevel{3}\enspace \R{解}{かい}の\R{個数}{こ|すう}を\R{判別式}{はん|べつ|しき}とグラフで\R{判断}{はん|だん}し、\R{場面}{ば|めん}に\R{合}{あ}う\R{解}{かい}を\R{選}{えら}べる。
\end{itemize}
\end{goal}

\begin{hatena}
\noindent\begin{minipage}[c]{0.6\linewidth}
3-9 で、\qty{20}{m/s} の\R{速}{はや}さで\R{真上}{ま|うえ}に\R{投}{な}げたボールの、$t$ \R{秒後}{びょう|ご}の\R{高}{たか}さは $y=20t-5t^2$（m）で、2\R{秒後}{びょう|ご}にいちばん\R{高}{たか}い \qty{20}{m} に\R{届}{とど}くことがわかりました。

では、\R{地面}{じ|めん}から \qty{15}{m} の\R{高}{たか}さにある\R{窓}{まど}の\R{前}{まえ}を、ボールが\R{通}{とお}るのは、\R{何秒後}{なん|びょう|ご}でしょう。

\R{答}{こた}えは、1つだけかな？
\end{minipage}\hfill
\begin{minipage}[c]{0.37\linewidth}\centering
\begin{tikzpicture}[x=0.78mm, y=0.78mm]
  % 1 m を 2 mm でかいた。窓は 15 m（30 mm）、最高点は 20 m（40 mm）
  \fill[bkGray!25] (14,0) rectangle (34,46);
  \draw[bkGray!80!black, line width=0.6pt] (14,0) rectangle (34,46);
  \foreach \yy in {6,18,30} {\foreach \xx in {17,26} {\fill[white] (\xx,\yy-3) rectangle ++(6,6);}}
  \fill[bkBlue!25] (17,27) rectangle ++(6,6);
  \draw[bkRed, line width=0.7pt, dash pattern=on 2pt off 1.5pt] (-2,30) -- (26,30);
  \fill[bkGreen!30] (-14,-2.5) rectangle (40,0);
  \draw[bkGreen!60!black, line width=0.8pt] (-14,0) -- (40,0);
  \draw[-{Stealth[length=3pt]}, bkOrange!80!black, line width=0.8pt] (-6,4) -- (-6,37);
  \draw[-{Stealth[length=3pt]}, bkOrange!80!black, line width=0.8pt] (2,37) -- (2,4);
  \draw[bkOrange!80!black, line width=0.8pt] (-6,37) arc (180:0:4);
  \bkthrkball{(-2,41)}{1.8}
  \node[font=\sffamily\bfseries\large, text=bkOrange!80!black, anchor=east] at (-7.5,30) {?};
  \node[font=\sffamily\bfseries\large, text=bkOrange!80!black, anchor=west] at (3.5,24) {?};
  \node[font=\sffamily\scriptsize, text=bkRed!80!black, anchor=south] at (8,30.8) {\qty{15}{m}};
\end{tikzpicture}
\end{minipage}
\end{hatena}

\clearpage
%% ---------------- やってみよう・思い出そう ----------------
\begin{yattemiyou}[\R{表}{ひょう}で、ボールの\R{高}{たか}さをさがそう]
$y=20t-5t^2$ に、$t=0$、0.5、1、…、4 を\R{入}{い}れて、\R{高}{たか}さを\R{表}{ひょう}にしました（\R{電卓}{でん|たく}で\R{確}{たし}かめよう）。
\begin{bkcenter}
\begin{bktable}{colspec={Q[l,m,wd=12mm] *{9}{Q[c,m,wd=9mm]}}, rowsep=1.5pt, colsep=2pt, cells={font=\small}}
$t$（s） & 0 & 0.5 & 1 & 1.5 & 2 & 2.5 & 3 & 3.5 & 4 \\
$y$（m） & 0 & 8.75 & 15 & 18.75 & 20 & 18.75 & 15 & 8.75 & 0 \\
\end{bktable}
\end{bkcenter}
\begin{enumerate}[label={\bkcircled[bkGreen]{\arabic*}}, leftmargin=2.4\zw, itemsep=1pt, topsep=3pt]
\item \R{高}{たか}さが \qty{15}{m} になるのは、\R{何秒後}{なん|びょう|ご}と\R{何秒後}{なん|びょう|ご}かな？
\item \R{高}{たか}さが \qty{20}{m} になるのは？　\qty{25}{m} になるのは？
\item \R{高}{たか}さが \qty{10}{m} になるのは、\R{何秒後}{なん|びょう|ご}かな？　\R{表}{ひょう}から\R{見}{み}つかるかな？
\end{enumerate}
\end{yattemiyou}

\noindent ①は、1\R{秒後}{びょう|ご}と3\R{秒後}{びょう|ご}です。1\R{秒後}{びょう|ご}は\R{上}{あ}がっていくとちゅう（\R{上}{のぼ}り）、3\R{秒後}{びょう|ご}は\R{下}{くだ}りてくるとちゅう（\R{下}{くだ}り）です。\R{答}{こた}えは2つありました。

②の \qty{20}{m} は、いちばん\R{高}{たか}い2\R{秒後}{びょう|ご}の1\R{回}{かい}だけです。\qty{25}{m} には、\R{一度}{いち|ど}も\R{届}{とど}きません。

③の \qty{10}{m} は、\R{表}{ひょう}にありません。0.5\R{秒}{びょう}（8.75 m）と1\R{秒}{びょう}（15 m）のあいだのどこかと、3\R{秒}{びょう}と3.5\R{秒}{びょう}のあいだのどこかです。\R{表}{ひょう}をどんなに\R{細}{ほそ}かくしても、ぴったりの\R{時刻}{じ|こく}は\R{見}{み}つからないかもしれません。

\R{表}{ひょう}でさがすのではなく、\R{式}{しき}を\R{計算}{けい|さん}して\R{時刻}{じ|こく}を\R{求}{もと}める\R{方法}{ほう|ほう}がほしくなりますね。それが、この\R{章}{しょう}で\R{学}{まな}ぶ「2\R{次方程式}{じ|ほう|てい|しき}を\R{解}{と}く」ことです。

\begin{omoidasou}
\bkrecall{2-8}{2\R{乗}{じょう}して9になる\R{数}{かず}は、3と $-3$ の2つ。$\sqrt{9}$ は、そのうちの\R{正}{せい}のほうの3。}
\bkrecall{2-16}{\R{方程式}{ほう|てい|しき}は、\R{両辺}{りょう|へん}に\R{同}{おな}じことをして（\R{移項}{い|こう}して）、$x=\cdots$ の\R{形}{かたち}にして\R{解}{と}く。}
\bkrecall{3-8}{$(x+a)(x+b)=x^2+(a+b)x+ab$、$(a+b)(a-b)=a^2-b^2$。\R{共通因数}{きょう|つう|いん|すう}でくくる：$3x^2-6x=3x(x-2)$。}
\bkrecall{3-8}{\R{平方完成}{へい|ほう|かん|せい}：$x^2-4x+3=(x-2)^2-1$（$-4$ の\R{半分}{はん|ぶん}の $-2$ を\R{使}{つか}う）。}
\bkrecall{3-9}{$y=-5t^2+20t=-5(t-2)^2+20$ は、\R{頂点}{ちょう|てん} $(2,\ 20)$ の、\R{下}{した}に\R{開}{ひら}いた\R{放物線}{ほう|ぶつ|せん}。}
\end{omoidasou}

\bkfukidashi{mina}{\R{表}{ひょう}にない\R{時刻}{じ|こく}も、\R{式}{しき}から\R{計算}{けい|さん}で\R{出}{だ}せたら、どんな\R{高}{たか}さでもわかるね。}

%% ---------------- 数学の道具 ----------------
\Needspace{14\baselineskip}
\bksection{doguu}{\R{数学}{すう|がく}の\R{道具}{どう|ぐ}}

\bksubsection{1}{$x^2$ のある\R{方程式}{ほう|てい|しき}：2\R{次方程式}{じ|ほう|てい|しき}}

\leavevmode\label{bk:3-10:niji}\bkhatenaref の\R{問}{と}いを、\R{式}{しき}にしてみましょう。\R{高}{たか}さ $20t-5t^2$ が15になる $t$ を\R{求}{もと}めるので、
\bkthrkline{$20t-5t^2=15$}
という\R{方程式}{ほう|てい|しき}になります。2-16 の\R{方程式}{ほう|てい|しき}とちがって、$t^2$（$t\times t$）の\R{項}{こう}があります。

この\R{式}{しき}を、あとで\R{解}{かい}きやすい\R{形}{かたち}に\R{直}{なお}しておきます。2-14 の\R{天}{てん}びんのように、\R{両辺}{りょう|へん}に\R{同}{おな}じことをしても、\R{等式}{とう|しき}はなりたったままです。
\bkthrkstep{1}{\R{右辺}{う|へん}を0にする}
\R{両辺}{りょう|へん}から15をひきます（\R{移項}{い|こう}、2-16）。\R{左辺}{さ|へん}は、\R{次数}{じ|すう}の\R{大}{おお}きい\R{順}{じゅん}にならべかえます。
\[ -5t^2+20t-15=0 \]
\bkthrkstep{2}{$t^2$ の\R{係数}{けい|すう}を1にする}
\R{両辺}{りょう|へん}を $-5$ でわります。\R{左辺}{さ|へん}は、\R{項}{こう}を1つずつ $-5$ でわります（2-12）。\R{右辺}{う|へん}は $0\div(-5)=0$ です。
\[ t^2-4t+3=0 \]
$-5t^2\div(-5)=t^2$、$20t\div(-5)=-4t$、$-15\div(-5)=3$ です。\R{符号}{ふ|ごう}が\R{全部}{ぜん|ぶ}\R{入}{い}れかわることに\R{注意}{ちゅう|い}しましょう。

\noindent\begin{minipage}[c]{0.56\linewidth}
やってみようで\R{見}{み}つけた $t=1$ と $t=3$ を、\R{左辺}{さ|へん}の $t^2-4t+3$ に\R{入}{い}れて\R{確}{たし}かめます（\R{右}{みぎ}の\R{表}{ひょう}）。$t=1$ なら $1-4+3=0$、$t=3$ なら $9-12+3=0$ で、どちらも\R{右辺}{う|へん}の0と\R{等}{ひと}しくなり、\R{方程式}{ほう|てい|しき}がなりたちます。$t=2$ では $4-8+3=-1$ で、なりたちません。
\end{minipage}\hfill
\begin{minipage}[c]{0.4\linewidth}\centering
\begin{bktable}{colspec={Q[c,m,wd=8mm] Q[c,m,wd=30mm]}, rowsep=1pt, cells={font=\small}}
$t$ & $t^2-4t+3$ \\
0 & $0-0+3=3$ \\
\SetCell{bg=bkYellow!30} 1 & \SetCell{bg=bkYellow!30} $1-4+3=0$ \\
2 & $4-8+3=-1$ \\
\SetCell{bg=bkYellow!30} 3 & \SetCell{bg=bkYellow!30} $9-12+3=0$ \\
4 & $16-16+3=3$ \\
\end{bktable}
\end{minipage}

\medskip
$t^2-4t+3=0$ のように、\R{文字}{も|じ}の2\R{乗}{じょう}の\R{項}{こう}があって、それより\R{大}{おお}きい\R{次数}{じ|すう}の\R{項}{こう}がない\R{方程式}{ほう|てい|しき}を、\textbf{2\R{次方程式}{じ|ほう|てい|しき}}といいます。\bkidx{にじほうていしき}{2\R{次方程式}{じ|ほう|てい|しき}}「2\R{次}{じ}」は、\R{文字}{も|じ}を2\R{回}{かい}かけた\R{項}{こう}がある、という\R{意味}{い|み}です（2-16 で「$t^2$ のある\R{方程式}{ほう|てい|しき}は2\R{次方程式}{じ|ほう|てい|しき}で、3-10 で\R{学}{まな}ぶ」と\R{予告}{よ|こく}していました）。

\R{方程式}{ほう|てい|しき}をなりたたせる\R{文字}{も|じ}の\R{値}{あたい}を、\R{方程式}{ほう|てい|しき}の\textbf{\R{解}{かい}}といいます。2\R{次方程式}{じ|ほう|てい|しき}には、\R{解}{かい}が2つあることがよくあります。\R{解}{かい}をすべて\R{求}{もと}めることを、\R{方程式}{ほう|てい|しき}を\textbf{\R{解}{と}く}といいます。\R{一般}{いっ|ぱん}に、2\R{次方程式}{じ|ほう|てい|しき}は
\bkthrkline{$ax^2+bx+c=0$\quad（$a$ は0でない\R{数}{かず}）}
の\R{形}{かたち}にまとめられます。$a$ が0だと $x^2$ の\R{項}{こう}が\R{消}{け}えて、1\R{次方程式}{じ|ほう|てい|しき}になってしまうからです。

\bksubsection{2}{2\R{乗}{じょう}して $k$ になる\R{数}{かず}：$x^2=k$}

\leavevmode\label{bk:3-10:root}いちばんかんたんな2\R{次方程式}{じ|ほう|てい|しき}から\R{始}{はじ}めましょう。ボールを\R{落}{お}とすと、$t$ \R{秒間}{びょう|かん}に\R{落}{お}ちる\R{距離}{きょ|り}は $5t^2$ m でした（2-13）。\qty{20}{m} \R{落}{お}ちるのは\R{何秒後}{なん|びょう|ご}でしょう。
\[ 5t^2=20 \]
\R{両辺}{りょう|へん}を5でわると、$t^2=4$ です。2\R{乗}{じょう}して4になる\R{数}{かず}をさがします。$2^2=4$ ですが、$(-2)^2=(-2)\times(-2)=4$ も4です（2-8）。だから、
\bkthrkline{$t=2$\quad または\quad $t=-2$}
の2つが\R{解}{かい}です。この2つを、まとめて $t=\pm2$ と\R{書}{か}きます。

\noindent\begin{minipage}[c]{0.52\linewidth}
\textbf{グラフで\R{見}{み}ると}\par\noindent $y=x^2$ のグラフ（3-9）と、\R{高}{たか}さ4の\R{横}{よこ}の\R{直線}{ちょく|せん} $y=4$ をかくと、2か\R{所}{ところ}で\R{交}{まじ}わります。\R{交}{まじ}わる\R{点}{てん}の $x$ が、$x^2=4$ の\R{解}{かい}の2と $-2$ です。\R{放物線}{ほう|ぶつ|せん}が\R{左右対称}{さ|ゆう|たい|しょう}なので、\R{交}{まじ}わる\R{点}{てん}も、0から\R{同}{おな}じ\R{距離}{きょ|り}の\R{左右}{さ|ゆう}に1つずつあります。
\end{minipage}\hfill
\begin{minipage}[c]{0.44\linewidth}\centering
\begin{tikzpicture}
\begin{axis}[bkplot, width=54mm, height=44mm, xmin=-3.2, xmax=3.6, ymin=-0.6, ymax=7.2, xtick={-2,2}, ytick=\empty,
  xlabel={$x$}, ylabel={$y$}]
  \addplot[bkcurve, domain=-2.6:2.6] {x^2};
  \addplot[bkcurve2, line width=1.2pt, domain=-3.1:3.1] {4};
  \addplot[bkGray, line width=0.6pt, dash pattern=on 2pt off 1.5pt] coordinates {(-2,0) (-2,4)};
  \addplot[bkGray, line width=0.6pt, dash pattern=on 2pt off 1.5pt] coordinates {(2,0) (2,4)};
  \addplot[only marks, mark=*, mark size=1.8pt, bkInk] coordinates {(-2,4) (2,4)};
  \node[font=\sffamily\scriptsize, text=bkOrange!80!black, anchor=south west] at (axis cs:2.7,4.1) {$y=4$};
\end{axis}
\end{tikzpicture}
\end{minipage}

\bkthrkhead{この「$\pm$」は、2-6 の「$\pm$」とちがう}
\noindent\begin{minipage}[c]{0.5\linewidth}
2-6 では、$\qty{50.0}{cm}\pm\qty{0.1}{cm}$ のように、$\pm$ を「49.9 から 50.1 までの、どこかにある」という\textbf{ずれのはば}に\R{使}{つか}いました。あいだの\R{数}{かず}も、ぜんぶふくみます。

ここの $t=\pm2$ は、「2と $-2$ の\textbf{ちょうど2つの\R{数}{かず}}」という\R{意味}{い|み}です。あいだの1や0は\R{解}{かい}ではありません。\R{同}{おな}じ\R{記号}{き|ごう}でも、\R{使}{つか}う\R{場面}{ば|めん}で\R{意味}{い|み}が\R{変}{か}わります。
\end{minipage}\hfill
\begin{minipage}[c]{0.46\linewidth}\centering
\begin{tikzpicture}[x=1mm, y=1mm]
  % 上：2-6 のずれのはば（帯）。下：x^2=4 の解（2つの点）
  \draw[bkInk, line width=0.7pt, -{Stealth[length=3pt]}] (0,18) -- (54,18);
  \fill[bkGreen!35] (17,16.6) rectangle (37,19.4);
  \foreach \xx/\lab in {17/49.9, 27/50.0, 37/50.1} {\draw[bkInk, line width=0.5pt] (\xx,16.8) -- (\xx,19.2); \node[font=\sffamily\scriptsize, anchor=north] at (\xx,15.6) {\lab};}
  \node[font=\sffamily\scriptsize, text=bkGreen!50!black, anchor=south] at (27,20.2) {2-6：あいだもぜんぶ};
  \draw[bkInk, line width=0.7pt, -{Stealth[length=3pt]}] (0,0) -- (54,0);
  \foreach \xx/\lab in {7/-2, 27/0, 47/2} {\draw[bkInk, line width=0.5pt] (\xx,-1.2) -- (\xx,1.2); \node[font=\sffamily\scriptsize, anchor=north] at (\xx,-2.4) {$\lab$};}
  \fill[bkRed] (7,0) circle (1.4); \fill[bkRed] (47,0) circle (1.4);
  \node[font=\sffamily\scriptsize, text=bkRed!80!black, anchor=south] at (27,1.8) {3-10：2つの\R{点}{てん}だけ};
\end{tikzpicture}
\end{minipage}

\medskip
\bkthrkhead{\R{場面}{ば|めん}に\R{合}{あ}う\R{解}{かい}をえらぶ}
ボールは、\R{落}{お}としはじめた $t=0$ より\R{後}{あと}の\R{時刻}{じ|こく}で\R{考}{かんが}えています。$t=-2$ は「2\R{秒前}{びょう|まえ}」なので、この\R{場面}{ば|めん}には\R{合}{あ}いません。\R{答}{こた}えは2\R{秒後}{びょう|ご}です。\R{方程式}{ほう|てい|しき}を\R{解}{と}いたら、\R{解}{かい}が\R{場面}{ば|めん}に\R{合}{あ}うかを、いつも\R{確}{たし}かめましょう。

\bkthrkhead{$k$ がいろいろな\R{数}{かず}のとき}
\begin{itemize}[leftmargin=1.6\zw, itemsep=2pt, topsep=1pt, label={\color{\bkTT}●}]
\item $x^2=7$：2\R{乗}{じょう}して7になる\R{整数}{せい|すう}はありません。2-8 の $\sqrt{\phantom{2}}$ を\R{使}{つか}って、$x=\pm\sqrt{7}$ と\R{書}{か}きます（$\sqrt{7}$ は\R{約}{やく}2.65）。
\item $x^2=0$：2\R{乗}{じょう}して0になる\R{数}{かず}は0だけなので、\R{解}{かい}は $x=0$ の1つです。
\item $x^2=-4$：\R{正}{せい}の\R{数}{かず}も\R{負}{ふ}の\R{数}{かず}も、2\R{乗}{じょう}するとプラスになります。だから、2\R{乗}{じょう}して $-4$ になる\R{数}{かず}は、この\R{巻}{かん}で\R{学}{まな}ぶ\R{数}{かず}の\R{中}{なか}にはありません。\R{解}{かい}はありません（このような\R{数}{かず}は 6-8 で\R{新}{あたら}しく\R{考}{かんが}えます）。
\end{itemize}

\bkthrkhead{かっこの2\R{乗}{じょう}のとき：$(x-2)^2=9$}
$x-2$ を、1つの\R{箱}{はこ}（2-11）と\R{見}{み}ます。\R{箱}{はこ}の\R{中身}{なか|み}を2\R{乗}{じょう}すると9になるので、\R{箱}{はこ}の\R{中身}{なか|み}は $3$ か $-3$ です。
\[ x-2=\pm3 \]
$x-2=3$ なら $x=5$、$x-2=-3$ なら $x=-1$ です（\R{両辺}{りょう|へん}に2をたす）。\R{確}{たし}かめると、$x=5$ のとき $(5-2)^2=3^2=9$、$x=-1$ のとき $(-1-2)^2=(-3)^2=9$ で、どちらもなりたちます。この「かっこの2\R{乗}{じょう}＝\R{数}{かず}」の\R{形}{かたち}は、あとで\R{平方完成}{へい|ほう|かん|せい}で\R{解}{と}くときに、もう\R{一度}{いち|ど}\R{使}{つか}います。

\begin{taisetsu}
$x^2=k$ の\R{解}{かい}（$k>0$）は、$x=\pm\sqrt{k}$（$\sqrt{k}$ と $-\sqrt{k}$ の2つ）。\par\vspace{2pt}
$k=0$ なら $x=0$ だけ。$k<0$ なら\R{解}{かい}はない。
\end{taisetsu}

\bksubsection{3}{かけて0なら、どちらかが0}

\leavevmode\label{bk:3-10:zero}\R{次}{つぎ}の □ と △ に\R{入}{はい}る\R{数}{かず}を\R{考}{かんが}えます。
\bkthrkline{□ × △ ＝ 0}
□ が3なら、$3\times△=0$ なので、△ は $0\div3=0$ です。□ が $-5$ でも、△ は0です。□ も △ も0でない\R{数}{かず}なら、かけ\R{算}{ざん}の\R{答}{こた}えは0になりません（2-2：\R{正}{せい}の\R{数}{かず}や\R{負}{ふ}の\R{数}{かず}どうしをかけても、0にはならない）。だから、かけて0になるのは、□ と △ の\textbf{\R{少}{すこ}なくともどちらか\R{一方}{いっ|ぽう}が0}のときだけです。

\noindent\begin{minipage}[c]{0.56\linewidth}
\textbf{\R{面積}{めん|せき}でも\R{考}{かんが}えられる}\par\noindent かけ\R{算}{ざん}は、\R{長方形}{ちょう|ほう|けい}の\R{面積}{めん|せき}でした（1-12）。\R{面積}{めん|せき}が0の\R{長方形}{ちょう|ほう|けい}は、たてか\R{横}{よこ}の\R{長}{なが}さが0で、ぺちゃんこにつぶれた\R{形}{かたち}です。たても\R{横}{よこ}も0でなければ、\R{面積}{めん|せき}は\R{必}{かなら}ず0より\R{大}{おお}きくなります。
\end{minipage}\hfill
\begin{minipage}[c]{0.4\linewidth}\centering
\begin{tikzpicture}[x=0.7mm, y=0.7mm]
  \filldraw[fill=bkBlue!20, draw=bkBlue!70!black, line width=0.7pt] (0,0) rectangle (16,12);
  \node[font=\sffamily\scriptsize, anchor=north] at (8,-1.5) {\R{面積}{めん|せき}は0でない};
  \filldraw[fill=bkBlue!20, draw=bkBlue!70!black, line width=0.7pt] (24,0) rectangle (40,4);
  \filldraw[fill=bkBlue!20, draw=bkBlue!70!black, line width=1.2pt] (46,0) -- (62,0);
  \draw[-{Stealth[length=3pt]}, bkInk, line width=0.6pt] (32,9) -- (32,5);
  \node[font=\sffamily\scriptsize, anchor=north] at (54,-1.5) {たて0で\R{面積}{めん|せき}0};
\end{tikzpicture}
\end{minipage}

\medskip
このきまりを\R{使}{つか}うと、かけ\R{算}{ざん}の\R{形}{かたち}をした\R{方程式}{ほう|てい|しき}が\R{解}{と}けます。
\bkthrkline{$x(x-2)=0$}
$x$ と $x-2$ をかけて0なので、$x=0$ か、$x-2=0$ のどちらかです。$x-2=0$ なら $x=2$ です。だから、\R{解}{かい}は $x=0$ と $x=2$ で、まとめて「$x=0,\ 2$」と\R{書}{か}きます。\R{確}{たし}かめ：$x=2$ なら $2\times(2-2)=2\times0=0$ です。

\bkthrkhead{0のときだけのきまり}
\noindent\begin{minipage}[c]{0.56\linewidth}
$(x-1)(x-3)=3$ を、「$x-1=3$ か $x-3=3$」として $x=4$、6 としてはいけません。$x=6$ を\R{入}{い}れると $5\times3=15$ で、3になりません。かけて3になる\R{組}{くみ}は、$1\times3$、$3\times1$、$0.5\times6$、… といくらでもあるので、1つにしぼれないのです。「どちらかが\R{決}{き}まる」のは、\R{答}{こた}えが0のときだけです。
\end{minipage}\hfill
\begin{minipage}[c]{0.4\linewidth}\centering
\begin{bktable}{colspec={Q[c,m,wd=12mm] Q[c,m,wd=12mm] Q[c,m,wd=16mm]}, rowsep=1pt, cells={font=\small}}
□ & △ & □×△ \\
1 & 3 & 3 \\
3 & 1 & 3 \\
0.5 & 6 & 3 \\
$-1$ & $-3$ & 3 \\
\end{bktable}
\end{minipage}

\begin{taisetsu}
$A\times B=0$ なら、$A=0$ または $B=0$。（\R{答}{こた}えが0のときだけ\R{使}{つか}える）
\end{taisetsu}

\bksubsection{4}{かけ\R{算}{ざん}の\R{形}{かたち}に\R{直}{なお}す：\R{因数分解}{いん|すう|ぶん|かい}}

\leavevmode\label{bk:3-10:insu}\noindent 3-8 で、かっこをはずす「\R{展開}{てん|かい}」を\R{学}{まな}びました。たとえば
\[ (t-1)(t-3)=t^2-3t-t+3=t^2-4t+3 \]
です。この\R{式}{しき}を、\R{右}{みぎ}から\R{左}{ひだり}へ\R{読}{よ}んでみましょう。$t^2-4t+3$ は $(t-1)(t-3)$ と、かけ\R{算}{ざん}の\R{形}{かたち}に\R{書}{か}けます。かけ\R{算}{ざん}の\R{形}{かたち}になれば、「かけて0なら、どちらかが0」（\pageref{bk:3-10:zero}ページ）が\R{使}{つか}えます。

$t^2-4t+3$ のような\R{和}{わ}の\R{形}{かたち}の\R{式}{しき}を、$(t-1)(t-3)$ のようなかけ\R{算}{ざん}の\R{形}{かたち}に\R{直}{なお}すことを、\textbf{\R{因数分解}{いん|すう|ぶん|かい}}といいます。\bkidx{いんすうぶんかい}{\R{因数分解}{いん|すう|ぶん|かい}}かけ\R{算}{ざん}を\R{作}{つく}っている1つ1つの\R{式}{しき}（$t-1$ や $t-3$）を、\textbf{\R{因数}{いん|すう}}といいます。$12=2\times6$ と、\R{数}{かず}をかけ\R{算}{ざん}に\R{分}{わ}けるのと\R{同}{おな}じ\R{考}{かんが}えです。

\begin{bkfigure}
\begin{tikzpicture}[x=1mm, y=1mm]
  \node[bk node, minimum width=36mm, minimum height=9mm, font=\small] (a) at (0,0) {$(t-1)(t-3)$};
  \node[bk node, minimum width=36mm, minimum height=9mm, font=\small] (b) at (70,0) {$t^2-4t+3$};
  \draw[bk link] ([yshift=2mm]a.east) -- ([yshift=2mm]b.west) node[bk link label, midway, above=0.6mm] {\R{展開}{てん|かい}（3-8）};
  \draw[bk link] ([yshift=-2mm]b.west) -- ([yshift=-2mm]a.east) node[bk link label, midway, below=0.6mm] {\R{因数分解}{いん|すう|ぶん|かい}};
\end{tikzpicture}
\bkfigcaption{\R{因数分解}{いん|すう|ぶん|かい}は、\R{展開}{てん|かい}の\R{逆}{ぎゃく}。}
\end{bkfigure}

\R{因数分解}{いん|すう|ぶん|かい}のしかたを、4つに\R{分}{わ}けて、1つずつ\R{見}{み}ていきましょう。

\bkthrkhead{その1：\R{共通因数}{きょう|つう|いん|すう}でくくる}
3-8 で\R{学}{まな}んだ「\R{共通因数}{きょう|つう|いん|すう}でくくる」も、\R{因数分解}{いん|すう|ぶん|かい}です。どの\R{項}{こう}にも\R{同}{おな}じ\R{文字}{も|じ}や\R{数}{かず}がふくまれていたら、それをかっこの\R{外}{そと}に\R{出}{だ}します（\R{分配法則}{ぶん|ぱい|ほう|そく}の\R{逆}{ぎゃく}、2-12）。
\begin{align*}
x^2-2x&=x\times x-2\times x=x(x-2)\\
3t^2-6t&=3t\times t-3t\times2=3t(t-2)\\
3x^2-3&=3\times x^2-3\times1=3(x^2-1)
\end{align*}
1\R{行}{ぎょう}めは、$x^2$ にも $2x$ にも $x$ がふくまれているので、$x$ でくくりました。2\R{行}{ぎょう}めは、$3t$ でくくりました。

\bkthrkhead{その2：\R{和}{わ}と\R{積}{せき}で、2つの\R{数}{かず}をさがす}
3-8 の\R{展開}{てん|かい}の\R{公式}{こう|しき}
\bkthrkline{$(x+a)(x+b)=x^2+(a+b)x+ab$}
を、\R{右}{みぎ}から\R{左}{ひだり}へ\R{読}{よ}みます。$x^2+\fbox{\R{和}{わ}}\,x+\fbox{\R{積}{せき}}$ の\R{形}{かたち}なら、\textbf{たして $x$ の\R{係数}{けい|すう}、かけて\R{定数項}{てい|すう|こう}}になる2つの\R{数}{かず} $a$、$b$ をさがせば、$(x+a)(x+b)$ と\R{因数分解}{いん|すう|ぶん|かい}できます。

\noindent\begin{minipage}[c]{0.5\linewidth}
\textbf{\R{面積図}{めん|せき|ず}で\R{見}{み}ると}\par\noindent $(x+2)(x+6)$ は、たて $x+2$、\R{横}{よこ} $x+6$ の\R{長方形}{ちょう|ほう|けい}の\R{面積}{めん|せき}です。\R{右}{みぎ}の\R{図}{ず}のように4つに\R{分}{わ}けると、$x^2$、$6x$、$2x$、12 の4つの\R{部分}{ぶ|ぶん}になり、たすと $x^2+8x+12$ です。$x$ の\R{係数}{けい|すう}の8は $2+6$（\R{和}{わ}）、12は $2\times6$（\R{積}{せき}）です。\R{反対}{はん|たい}に、$x^2+8x+12$ を\R{長方形}{ちょう|ほう|けい}にならべるには、たして8、かけて12の2つの\R{数}{かず}を\R{見}{み}つければよいのです。
\end{minipage}\hfill
\begin{minipage}[c]{0.46\linewidth}\centering
\begin{tikzpicture}[x=1mm, y=1mm]
  \bkthrktile{(0,0)}{24}{24}{bkBlue}{$x^2$}
  \bkthrktile{(24,0)}{18}{24}{bkGreen}{$6x$}
  \bkthrktile{(0,24)}{24}{8}{bkOrange}{$2x$}
  \bkthrktile{(24,24)}{18}{8}{bkRed}{12}
  \draw[decorate, decoration={brace, amplitude=2.5pt}, line width=0.6pt] (0,33.5) -- (24,33.5) node[midway, above=1.5mm, font=\footnotesize] {$x$};
  \draw[decorate, decoration={brace, amplitude=2.5pt}, line width=0.6pt] (24,33.5) -- (42,33.5) node[midway, above=1.5mm, font=\footnotesize] {6};
  \draw[decorate, decoration={brace, amplitude=2.5pt}, line width=0.6pt] (-1.5,0) -- (-1.5,24) node[midway, left=1.5mm, font=\footnotesize] {$x$};
  \draw[decorate, decoration={brace, amplitude=2.5pt}, line width=0.6pt] (-1.5,24) -- (-1.5,32) node[midway, left=1.5mm, font=\footnotesize] {2};
\end{tikzpicture}
\end{minipage}

\medskip
\textbf{\R{例}{れい}　$x^2-8x+12$}\quad たして $-8$、かけて12 になる2つの\R{数}{かず}をさがします。かけて12になる\R{整数}{せい|すう}の\R{組}{くみ}を、ぜんぶ\R{書}{か}き\R{出}{だ}して、たした\R{数}{かず}を\R{調}{しら}べます。
\begin{bkcenter}
\begin{bktable}{colspec={Q[l,m,wd=20mm] *{6}{Q[c,m,wd=14mm]}}, rowsep=1.5pt, colsep=3pt, cells={font=\small}}
かけて12の\R{組}{くみ} & 1と12 & 2と6 & 3と4 & $-1$と$-12$ & $-2$と$-6$ & $-3$と$-4$ \\
たした\R{数}{かず} & 13 & 8 & 7 & $-13$ & \SetCell{bg=bkYellow!30} $-8$ & $-7$ \\
\end{bktable}
\bktabcaption{\R{和}{わ}が $-8$、\R{積}{せき}が12になるのは、$-2$ と $-6$。}
\end{bkcenter}
\R{和}{わ}が $-8$、\R{積}{せき}が12 になる2つの\R{数}{かず}は $-2$ と $-6$ なので、
\[ x^2-8x+12=(x-2)(x-6) \]
です。\R{確}{たし}かめに\R{展開}{てん|かい}すると、$x^2-6x-2x+12=x^2-8x+12$ で、もとにもどります。

\textbf{\R{例}{れい}　$t^2+2t-15$}\quad かけて $-15$、たして2 になる2つの\R{数}{かず}をさがします。かけてマイナスになるのは、プラスとマイナスの\R{組}{くみ}です。
\begin{bkcenter}
\begin{bktable}{colspec={Q[l,m,wd=26mm] *{4}{Q[c,m,wd=16mm]}}, rowsep=1.5pt}
かけて $-15$ の\R{組}{くみ} & 1と$-15$ & $-1$と15 & 3と$-5$ & $-3$と5 \\
たした\R{数}{かず} & $-14$ & 14 & $-2$ & \SetCell{bg=bkYellow!30} 2 \\
\end{bktable}
\end{bkcenter}
5と $-3$ なので、$t^2+2t-15=(t+5)(t-3)$ です。\R{確}{たし}かめに\R{展開}{てん|かい}すると、
\[ (t+5)(t-3)=t^2-3t+5t-15=t^2+2t-15 \]
で、もとにもどります。

\begin{motto}[2]
\textbf{\R{符号}{ふ|ごう}のヒント}\par
かけた\R{数}{かず}（\R{積}{せき}）がプラスなら、2つの\R{数}{かず}は\R{同}{おな}じ\R{符号}{ふ|ごう}です。そのとき、たした\R{数}{かず}（\R{和}{わ}）がプラスなら2つともプラス、マイナスなら2つともマイナスです。\R{積}{せき}がマイナスなら、2つの\R{数}{かず}はちがう\R{符号}{ふ|ごう}で、\R{大}{おお}きさの\R{大}{おお}きいほうが、\R{和}{わ}と\R{同}{おな}じ\R{符号}{ふ|ごう}になります。ヒントでしぼってから、\R{表}{ひょう}で\R{確}{たし}かめると\R{早}{はや}く\R{見}{み}つかります。
\end{motto}

\bkthrkhead{その3：$x^2-a^2=(x+a)(x-a)$}
3-8 の $(a+b)(a-b)=a^2-b^2$ を\R{右}{みぎ}から\R{左}{ひだり}へ\R{読}{よ}むと、「2\R{乗}{じょう}ひく2\R{乗}{じょう}」は、\R{和}{わ}と\R{差}{さ}のかけ\R{算}{ざん}に\R{因数分解}{いん|すう|ぶん|かい}できます。
\[ x^2-1=x^2-1^2=(x+1)(x-1),\qquad x^2-9=x^2-3^2=(x+3)(x-3) \]
\noindent\begin{minipage}[c]{0.56\linewidth}
\textbf{\R{和}{わ}と\R{積}{せき}でも、\R{同}{おな}じ\R{答}{こた}えになる}\par\noindent $x^2-1$ は、$x^2+0x-1$ と\R{見}{み}ることができます。たして0、かけて $-1$ の2つの\R{数}{かず}は、1と $-1$ です。だから $(x+1)(x-1)$。その2の\R{方法}{ほう|ほう}でも、\R{同}{おな}じ\R{答}{こた}えが\R{出}{で}ます。

\textbf{\R{面積図}{めん|せき|ず}で\R{見}{み}ると}\par\noindent 1\R{辺}{ぺん} $x$ の\R{正方形}{せい|ほう|けい}から、1\R{辺}{ぺん}1の\R{正方形}{せい|ほう|けい}を\R{切}{き}りとると（\R{右}{みぎ}の\R{図}{ず}の\R{左}{ひだり}）、\R{残}{のこ}りは $x^2-1$ です。\R{下}{した}の\R{細}{ほそ}い\R{部分}{ぶ|ぶん}（たて1、\R{横}{よこ} $x-1$）を、\R{右}{みぎ}のはしに\R{立}{た}ててつなぐと、たて $x-1$、\R{横}{よこ} $x+1$ の\R{長方形}{ちょう|ほう|けい}になります。
\end{minipage}\hfill
\begin{minipage}[c]{0.4\linewidth}\centering
\begin{tikzpicture}[x=0.7mm, y=0.7mm]
  % x=20 mm、1=5 mm
  \filldraw[fill=bkBlue!22, draw=bkBlue!75!black, line width=0.6pt] (0,5) rectangle (15,20);
  \filldraw[fill=bkGreen!22, draw=bkGreen!75!black, line width=0.6pt] (0,0) rectangle (15,5);
  \filldraw[fill=bkBlue!22, draw=bkBlue!75!black, line width=0.6pt] (15,5) rectangle (20,20);
  \draw[bkGray, line width=0.6pt, dash pattern=on 1.5pt off 1pt] (15,0) rectangle (20,5);
  \draw[-{Stealth[length=3pt]}, bkInk, line width=0.7pt] (23,10) -- (29,10);
  \filldraw[fill=bkBlue!22, draw=bkBlue!75!black, line width=0.6pt] (32,5) rectangle (52,20);
  \filldraw[fill=bkGreen!22, draw=bkGreen!75!black, line width=0.6pt] (52,5) rectangle (57,20);
  \draw[decorate, decoration={brace, amplitude=2.5pt}, line width=0.6pt] (32,21.5) -- (57,21.5) node[midway, above=1.5mm, font=\footnotesize] {$x+1$};
  \draw[decorate, decoration={brace, amplitude=2.5pt, mirror}, line width=0.6pt] (58.5,5) -- (58.5,20) node[midway, right=1.5mm, font=\footnotesize] {$x-1$};
  \draw[decorate, decoration={brace, amplitude=2.5pt}, line width=0.6pt] (0,21.5) -- (20,21.5) node[midway, above=1.5mm, font=\footnotesize] {$x$};
\end{tikzpicture}
\end{minipage}

\bkthrkhead{その4：くくってから、さらに\R{因数分解}{いん|すう|ぶん|かい}}
くくったあと、かっこの\R{中}{なか}がさらに\R{因数分解}{いん|すう|ぶん|かい}できることがあります。
\begin{align*}
3x^2-3&=3(x^2-1)=3(x+1)(x-1)\\
12x^2-96x+144&=12(x^2-8x+12)=12(x-2)(x-6)
\end{align*}
2\R{行}{ぎょう}めは、どの\R{項}{こう}も12でわれるので、\R{先}{さき}に12でくくりました。かっこの\R{中}{なか}は、\R{和}{わ}が $-8$、\R{積}{せき}が12 の\R{形}{かたち}です。\R{大}{おお}きな\R{数}{かず}のままさがすより、ずっと\R{楽}{らく}になります。

\begin{taisetsu}
\R{因数分解}{いん|すう|ぶん|かい}の\R{手順}{て|じゅん}\par
①\enspace どの\R{項}{こう}にも\R{共通}{きょう|つう}な\R{数}{かず}や\R{文字}{も|じ}があれば、くくる。\par
②\enspace 「2\R{乗}{じょう}ひく2\R{乗}{じょう}」なら、$x^2-a^2=(x+a)(x-a)$。\par
③\enspace $x^2+\fbox{\R{和}{わ}}\,x+\fbox{\R{積}{せき}}$ なら、たして\R{和}{わ}・かけて\R{積}{せき}の2つの\R{数}{かず}をさがす。\par
④\enspace \R{展開}{てん|かい}して、もとの\R{式}{しき}にもどるか\R{確}{たし}かめる。
\end{taisetsu}

\bksubsection{5}{\R{因数分解}{いん|すう|ぶん|かい}で2\R{次方程式}{じ|ほう|てい|しき}を\R{解}{と}く}

\leavevmode\label{bk:3-10:toku}\R{道具}{どう|ぐ}がそろいました。はてな？の\R{方程式}{ほう|てい|しき}を\R{解}{かい}きましょう。
\bkthrkstep{1}{\R{右辺}{う|へん}を0にして、\R{整理}{せい|り}する}
\pageref{bk:3-10:niji}ページと\R{同}{おな}じように、\R{移項}{い|こう}して\R{両辺}{りょう|へん}を $-5$ でわります。
\[ 20t-5t^2=15\qquad\text{より}\qquad t^2-4t+3=0 \]
\bkthrkstep{2}{\R{左辺}{さ|へん}を\R{因数分解}{いん|すう|ぶん|かい}する}
たして $-4$、かけて3 の2つの\R{数}{かず}は $-1$ と $-3$ です。
\[ (t-1)(t-3)=0 \]
\bkthrkstep{3}{かけて0なら、どちらかが0}
$t-1=0$ か $t-3=0$ なので、$t=1,\ 3$ です。
\bkthrkstep{4}{\R{場面}{ば|めん}に\R{合}{あ}うかを\R{確}{たし}かめる}
どちらも0より\R{後}{あと}の\R{時刻}{じ|こく}で、ボールが\R{地面}{じ|めん}にもどる4\R{秒後}{びょう|ご}より\R{前}{まえ}なので、\R{場面}{ば|めん}に\R{合}{あ}います。

\noindent\begin{minipage}[c]{0.5\linewidth}\centering
\begin{tikzpicture}
\begin{axis}[bkplot, width=60mm, height=50mm, xmin=0, xmax=4.8, ymin=0, ymax=24, xtick={1,2,3,4}, ytick={5,10,15,20},
  xlabel={$t$（s）}, ylabel={$y$（m）}]
  \addplot[bkcurve, domain=0:4] {20*x-5*x^2};
  \addplot[bkcurve2, line width=1.1pt, domain=0:4.4] {15};
  \addplot[bkGray, line width=0.6pt, dash pattern=on 2pt off 1.5pt] coordinates {(1,0) (1,15)};
  \addplot[bkGray, line width=0.6pt, dash pattern=on 2pt off 1.5pt] coordinates {(3,0) (3,15)};
  \addplot[only marks, mark=*, mark size=1.9pt, bkInk] coordinates {(1,15) (3,15)};
  \node[font=\sffamily\scriptsize, text=bkInk, anchor=east] at (axis cs:0.95,17.3) {\R{上}{のぼ}り};
  \node[font=\sffamily\scriptsize, text=bkInk, anchor=west] at (axis cs:3.05,17.3) {\R{下}{くだ}り};
\end{axis}
\end{tikzpicture}
\end{minipage}\hfill
\begin{minipage}[c]{0.46\linewidth}
\bkhatenaref の\R{答}{こた}え：ボールが\R{高}{たか}さ \qty{15}{m} の\R{窓}{まど}の\R{前}{まえ}を\R{通}{とお}るのは、\textbf{1\R{秒後}{びょう|ご}と3\R{秒後}{びょう|ご}}です。1\R{秒後}{びょう|ご}は\R{上}{のぼ}りで、3\R{秒後}{びょう|ご}は\R{下}{くだ}りです。

グラフでは、\R{放物線}{ほう|ぶつ|せん}と、\R{高}{たか}さ15の\R{直線}{ちょく|せん}（オレンジ）が2か\R{所}{ところ}で\R{交}{まじ}わります。\R{放物線}{ほう|ぶつ|せん}は\R{頂点}{ちょう|てん}（2\R{秒}{びょう}）をはさんで\R{左右対称}{さ|ゆう|たい|しょう}なので、2つの\R{時刻}{じ|こく}も2\R{秒}{びょう}から\R{同}{おな}じ1\R{秒}{びょう}ずつはなれています。\R{答}{こた}えが2つ\R{出}{で}るのは、\R{上}{のぼ}りと\R{下}{くだ}りがあるからです。
\end{minipage}

\begin{taisetsu}
2\R{次方程式}{じ|ほう|てい|しき}を\R{因数分解}{いん|すう|ぶん|かい}で\R{解}{と}く\R{手順}{て|じゅん}\par
①\enspace \R{移項}{い|こう}して、\R{右辺}{う|へん}を0にする。\quad ②\enspace \R{左辺}{さ|へん}を\R{因数分解}{いん|すう|ぶん|かい}する。\par
③\enspace かけて0なら、どちらかが0。\quad ④\enspace \R{場面}{ば|めん}に\R{合}{あ}う\R{解}{と}か\R{確}{たし}かめる。
\end{taisetsu}

\bksubsection{6}{\R{因数分解}{いん|すう|ぶん|かい}できないときは、\R{平方完成}{へい|ほう|かん|せい}で\R{解}{と}く}

\leavevmode\label{bk:3-10:heihou}やってみようの③、「\R{高}{たか}さ \qty{10}{m} になるのは\R{何秒後}{なん|びょう|ご}？」にもどります。
\[ 20t-5t^2=10\quad\text{より}\quad t^2-4t+2=0 \]
たして $-4$、かけて2 になる\R{整数}{せい|すう}の\R{組}{くみ}をさがすと、1と2（\R{和}{わ}3）、$-1$ と $-2$（\R{和}{わ} $-3$）しかなく、\R{見}{み}つかりません。\R{整数}{せい|すう}ではうまく\R{因数分解}{いん|すう|ぶん|かい}できないのです。

そこで、3-8・3-9 の\R{平方完成}{へい|ほう|かん|せい}を\R{使}{つか}います。「（かっこ）の2\R{乗}{じょう}＝\R{数}{かず}」の\R{形}{かたち}にすれば、$x^2=k$（\pageref{bk:3-10:root}ページ）と\R{同}{おな}じように\R{解}{と}けるからです。
\bkthrkstep{1}{\R{平方完成}{へい|ほう|かん|せい}する}
$t$ の\R{係数}{けい|すう} $-4$ の\R{半分}{はん|ぶん}は $-2$。$(t-2)^2=t^2-4t+4$ なので、4をたして4をひきます。
\[ t^2-4t+2=t^2-4t+4-4+2=(t-2)^2-2 \]
\bkthrkstep{2}{「かっこの2\R{乗}{じょう}＝\R{数}{かず}」の\R{形}{かたち}にする}
\R{方程式}{ほう|てい|しき}は $(t-2)^2-2=0$ になります。\R{両辺}{りょう|へん}に2をたして、
\[ (t-2)^2=2 \]
\bkthrkstep{3}{$\pm\sqrt{\phantom{2}}$ をとる}
$t-2$ は、2\R{乗}{じょう}して2になる\R{数}{かず}なので、$\sqrt{2}$ か $-\sqrt{2}$ です。
\[ t-2=\pm\sqrt{2}\qquad\text{だから}\qquad t=2\pm\sqrt{2} \]
\bkthrkstep{4}{\R{数}{かず}の\R{大}{おお}きさを\R{確}{たし}かめる}
$\sqrt{2}$ は\R{約}{やく}1.414（2-8）なので、$2-\sqrt{2}$ は\R{約}{やく}0.59、$2+\sqrt{2}$ は\R{約}{やく}3.41 です。やってみようの\R{表}{ひょう}で\R{見}{み}た「0.5と1のあいだ」「3と3.5のあいだ」に、ぴったり\R{合}{あ}います。\R{電卓}{でん|たく}で $t=0.586$ を\R{入}{い}れると、$20\times0.586-5\times0.586^2$ は\R{約}{やく}10.0 です。

ボールは、\R{約}{やく}0.59\R{秒後}{びょう|ご}（\R{上}{のぼ}り）と\R{約}{やく}3.41\R{秒後}{びょう|ご}（\R{下}{くだ}り）に、\R{高}{たか}さ \qty{10}{m} を\R{通}{とお}ります。\R{平方完成}{へい|ほう|かん|せい}を\R{使}{つか}えば、\R{因数分解}{いん|すう|ぶん|かい}できない2\R{次方程式}{じ|ほう|てい|しき}も\R{解}{と}けるのです。

\bksubsection{7}{$\sqrt{\phantom{2}}$ のついた\R{数}{かず}の\R{計算}{けい|さん}}

\leavevmode\label{bk:3-10:sqrt}\R{平方完成}{へい|ほう|かん|せい}や、\R{次}{つぎ}の\R{解}{かい}の\R{公式}{こう|しき}で\R{解}{と}くと、\R{答}{こた}えに $\sqrt{\phantom{2}}$ がよく\R{出}{で}てきます。$\sqrt{\phantom{2}}$ の\R{数}{かず}を、かんたんな\R{形}{かたち}に\R{直}{なお}す\R{方法}{ほう|ほう}を2つ\R{学}{まな}びましょう。

\bkthrkhead{$\sqrt{\phantom{2}}$ の\R{中}{なか}を\R{小}{ちい}さくする：$\sqrt{12}=2\sqrt{3}$}
2-8 で、$\sqrt{2}\times\sqrt{3}=\sqrt{6}$ のように、$\sqrt{a}\times\sqrt{b}=\sqrt{ab}$ がなりたつことを\R{確}{たし}かめました。これを\R{反対}{はん|たい}に\R{使}{つか}います。12を「2\R{乗}{じょう}の\R{数}{かず}×ほかの\R{数}{かず}」に\R{分}{わ}けると、$12=4\times3=2^2\times3$ です。
\[ \sqrt{12}=\sqrt{4\times3}=\sqrt{4}\times\sqrt{3}=2\times\sqrt{3}=2\sqrt{3} \]
\R{電卓}{でん|たく}で\R{確}{たし}かめると、$\sqrt{12}=3.4641\cdots$、$2\times\sqrt{3}=2\times1.7320\cdots=3.4641\cdots$ で、\R{同}{おな}じです。$\sqrt{\phantom{2}}$ の\R{中}{なか}は、できるだけ\R{小}{ちい}さい\R{整数}{せい|すう}にしておくと、\R{大}{おお}きさが\R{見}{み}やすく、\R{計算}{けい|さん}もしやすくなります。

\R{中}{なか}の\R{数}{かず}を、$4=2^2$、$9=3^2$、$16=4^2$、$25=5^2$ … でわれるか\R{調}{しら}べるのがコツです。$\sqrt{18}=\sqrt{9\times2}=3\sqrt{2}$、$\sqrt{20}=\sqrt{4\times5}=2\sqrt{5}$、$\sqrt{50}=\sqrt{25\times2}=5\sqrt{2}$ です。

\bkthrkhead{\R{分母}{ぶん|ぼ}の $\sqrt{\phantom{2}}$ をなくす：$\dfrac{1}{\sqrt{2}}=\dfrac{\sqrt{2}}{2}$}
$\dfrac{1}{\sqrt{2}}$ は、$1\div1.41421\cdots$ で、\R{手}{て}で\R{計算}{けい|さん}するのは\R{大変}{たい|へん}です。そこで、\R{分母}{ぶん|ぼ}と\R{分子}{ぶん|し}に\R{同}{おな}じ $\sqrt{2}$ をかけます。\R{分母}{ぶん|ぼ}と\R{分子}{ぶん|し}に\R{同}{おな}じ\R{数}{かず}をかけても、\R{分数}{ぶん|すう}の\R{大}{おお}きさは\R{変}{か}わりません（2-4）。
\[ \frac{1}{\sqrt{2}}=\frac{1\times\sqrt{2}}{\sqrt{2}\times\sqrt{2}}=\frac{\sqrt{2}}{2} \]
\R{分母}{ぶん|ぼ}の $\sqrt{2}\times\sqrt{2}$ が2になるのは、$\sqrt{2}$ が「2\R{乗}{じょう}すると2になる\R{数}{かず}」だからです。これなら、
\[ \frac{\sqrt{2}}{2}=1.41421\cdots\div2=0.70710\cdots \]
と、かんたんに\R{計算}{けい|さん}できます。\R{分母}{ぶん|ぼ}の $\sqrt{\phantom{2}}$ をなくすことを、\R{分母}{ぶん|ぼ}の\textbf{\R{有理化}{ゆう|り|か}}といいます。\bkidx{ゆうりか}{\R{有理化}{ゆう|り|か}}（$\sqrt{2}$ のような\R{数}{かず}は\R{無理数}{む|り|すう}、\R{整数}{せい|すう}や\R{分数}{ぶん|すう}は\R{有理数}{ゆう|り|すう}でした。2-9）

\bksubsection{8}{どんな2\R{次方程式}{じ|ほう|てい|しき}も\R{解}{と}ける：\R{解}{かい}の\R{公式}{こう|しき}}

\leavevmode\label{bk:3-10:kousiki}\pageref{bk:3-10:heihou}ページの\R{平方完成}{へい|ほう|かん|せい}の\R{手順}{て|じゅん}は、どんな2\R{次方程式}{じ|ほう|てい|しき} $ax^2+bx+c=0$ にも\R{使}{つか}えます。\R{数}{かず}の\R{例}{れい}（\R{左}{ひだり}）と、\R{文字}{も|じ}のまま（\R{右}{みぎ}）とを、ならべて\R{計算}{けい|さん}してみましょう。\R{左}{ひだり}の\R{例}{れい}は $2x^2+3x-1=0$ で、$a=2$、$b=3$、$c=-1$ です。

\begin{bkcenter}
\begin{bktable}{colspec={Q[l,m,wd=34mm] Q[c,m,wd=38mm] Q[c,m,wd=46mm]}, rowsep=3pt, cells={font=\small}, row{1}={font=\sffamily\bfseries\small}}
すること & $2x^2+3x-1=0$ & $ax^2+bx+c=0$ \\
① \R{両辺}{りょう|へん}を $a$ でわる & $x^2+\dfrac{3}{2}x-\dfrac{1}{2}=0$ & $x^2+\dfrac{b}{a}x+\dfrac{c}{a}=0$ \\
② \R{数}{かず}の\R{項}{こう}を\R{右辺}{う|へん}へ & $x^2+\dfrac{3}{2}x=\dfrac{1}{2}$ & $x^2+\dfrac{b}{a}x=-\dfrac{c}{a}$ \\
③ $x$ の\R{係数}{けい|すう}の\R{半分}{はん|ぶん}の2\R{乗}{じょう}を、\R{両辺}{りょう|へん}にたす & $+\left(\dfrac{3}{4}\right)^2=\dfrac{9}{16}$ & $+\left(\dfrac{b}{2a}\right)^2=\dfrac{b^2}{4a^2}$ \\
\end{bktable}
\end{bkcenter}
③で、\R{左辺}{さ|へん}は「（かっこ）の2\R{乗}{じょう}」の\R{形}{かたち}にまとまる\R{準備}{じゅん|び}ができました。\R{続}{つづ}けます。
\begin{bkcenter}
\begin{bktable}{colspec={Q[l,m,wd=34mm] Q[c,m,wd=38mm] Q[c,m,wd=46mm]}, rowsep=3pt, cells={font=\small}, row{1}={font=\sffamily\bfseries\small}}
すること & $2x^2+3x-1=0$ & $ax^2+bx+c=0$ \\
④ \R{左辺}{さ|へん}を（かっこ）の2\R{乗}{じょう}にする & $\left(x+\dfrac{3}{4}\right)^2=\dfrac{17}{16}$ & $\left(x+\dfrac{b}{2a}\right)^2=\dfrac{b^2-4ac}{4a^2}$ \\
⑤ $\pm\sqrt{\phantom{2}}$ をとる & $x+\dfrac{3}{4}=\pm\dfrac{\sqrt{17}}{4}$ & $x+\dfrac{b}{2a}=\pm\dfrac{\sqrt{b^2-4ac}}{2a}$ \\
⑥ \R{移項}{い|こう}する & $x=\dfrac{-3\pm\sqrt{17}}{4}$ & $x=\dfrac{-b\pm\sqrt{b^2-4ac}}{2a}$ \\
\end{bktable}
\bktabcaption{\R{平方完成}{へい|ほう|かん|せい}の\R{手順}{て|じゅん}を、\R{文字}{も|じ}のまま\R{進}{すす}める。}
\end{bkcenter}

1\R{行}{ぎょう}ずつ、なぜそうなるかを\R{確}{たし}かめます。
\begin{itemize}[leftmargin=1.6\zw, itemsep=2pt, topsep=1pt, label={\color{\bkTT}●}]
\item ①：$x^2$ の\R{係数}{けい|すう}を1にするためです。$a$ は0でないので、わってかまいません。
\item ③：$x$ の\R{係数}{けい|すう} $\dfrac{3}{2}$ の\R{半分}{はん|ぶん}は $\dfrac{3}{4}$、$\dfrac{b}{a}$ の\R{半分}{はん|ぶん}は $\dfrac{b}{2a}$ です（3-8 の\R{平方完成}{へい|ほう|かん|せい}と\R{同}{おな}じ）。
\item ④の\R{右辺}{う|へん}：\R{左}{ひだり}は $\dfrac{1}{2}+\dfrac{9}{16}=\dfrac{8}{16}+\dfrac{9}{16}=\dfrac{17}{16}$。\R{右}{みぎ}は、$-\dfrac{c}{a}=-\dfrac{4ac}{4a^2}$ と\R{分母}{ぶん|ぼ}をそろえて（\R{通分}{つう|ぶん}、2-4）、$\dfrac{b^2}{4a^2}-\dfrac{4ac}{4a^2}=\dfrac{b^2-4ac}{4a^2}$。
\item ⑤：\R{分母}{ぶん|ぼ}の $\sqrt{16}=4$、$\sqrt{4a^2}=2a$ です（$a$ がマイナスのときも、$\pm$ が\R{両方}{りょう|ほう}の\R{符号}{ふ|ごう}をふくむので、\R{同}{おな}じ\R{式}{しき}になります）。
\end{itemize}

こうしてできた\R{式}{しき}を、2\R{次方程式}{じ|ほう|てい|しき}の\textbf{\R{解}{かい}の\R{公式}{こう|しき}}といいます。\bkidx{かいのこうしき}{\R{解}{かい}の\R{公式}{こう|しき}}
\begin{taisetsu}
\textbf{\R{解}{かい}の\R{公式}{こう|しき}}\par
$ax^2+bx+c=0$ の\R{解}{かい}は\quad $x=\dfrac{-b\pm\sqrt{b^2-4ac}}{2a}$
\end{taisetsu}
\R{覚}{おぼ}えていなくても、\R{平方完成}{へい|ほう|かん|せい}の\R{手順}{て|じゅん}をたどれば、いつでも\R{作}{つく}り\R{直}{なお}せます。

\bkthrkhead{はてな？の\R{方程式}{ほう|てい|しき}でも\R{確}{たし}かめる}
$t^2-4t+3=0$ は、$a=1$、$b=-4$、$c=3$ です。$-b=-(-4)=4$、$b^2=(-4)^2=16$ に\R{注意}{ちゅう|い}して\R{入}{い}れると、
\[ t=\frac{4\pm\sqrt{16-4\times1\times3}}{2\times1}=\frac{4\pm\sqrt{4}}{2}=\frac{4\pm2}{2} \]
$+$ のほうが $\dfrac{6}{2}=3$、$-$ のほうが $\dfrac{2}{2}=1$ で、\R{因数分解}{いん|すう|ぶん|かい}で\R{求}{もと}めた1と3に\R{合}{あ}います。\R{因数分解}{いん|すう|ぶん|かい}できるときは\R{因数分解}{いん|すう|ぶん|かい}のほうが\R{早}{はや}く、できないときは\R{解}{かい}の\R{公式}{こう|しき}を\R{使}{つか}います。

\bksubsection{9}{\R{解}{かい}が2つ・1つ・ない：\R{判別式}{はん|べつ|しき}}

\leavevmode\label{bk:3-10:hanbetsu}やってみようの②を、\R{式}{しき}で\R{考}{かんが}えます。\R{高}{たか}さ15 m、20 m、25 m になる\R{時刻}{じ|こく}の\R{方程式}{ほう|てい|しき}は、\R{次}{つぎ}のようになります（どれも、\R{両辺}{りょう|へん}から\R{高}{たか}さをひいて $-5$ でわった）。

\begin{bkcenter}
\begin{bktable}{colspec={Q[c,m,wd=10mm] Q[c,m,wd=27mm] Q[c,m,wd=26mm] Q[c,m,wd=22mm] Q[c,m,wd=22mm]}, colsep=3pt, rowsep=2pt, cells={font=\small}, row{1}={font=\sffamily\bfseries\small}}
\R{高}{たか}さ & \R{方程式}{ほう|てい|しき} & \R{平方完成}{へい|ほう|かん|せい} & $b^2-4ac$ & \R{解}{かい} \\
\qty{15}{m} & $t^2-4t+3=0$ & $(t-2)^2=1$ & $16-12=4$ & 2つ（1、3） \\
\qty{20}{m} & $t^2-4t+4=0$ & $(t-2)^2=0$ & $16-16=0$ & 1つ（2） \\
\qty{25}{m} & $t^2-4t+5=0$ & $(t-2)^2=-1$ & $16-20=-4$ & ない \\
\end{bktable}
\end{bkcenter}

\noindent\begin{minipage}[c]{0.5\linewidth}\centering
\begin{tikzpicture}
\begin{axis}[bkplot, width=60mm, height=52mm, xmin=0, xmax=5.2, ymin=0, ymax=28, xtick={1,2,3,4}, ytick={15,20,25},
  xlabel={$t$（s）}, ylabel={$y$（m）}]
  \addplot[bkcurve, domain=0:4] {20*x-5*x^2};
  \addplot[bkGreen!70!black, line width=1.1pt, domain=0:4.4] {15};
  \addplot[bkOrange, line width=1.1pt, domain=0:4.4] {20};
  \addplot[bkRed, line width=1.1pt, domain=0:4.4] {25};
  \addplot[only marks, mark=*, mark size=1.8pt, bkInk] coordinates {(1,15) (3,15) (2,20)};
  \node[font=\sffamily\scriptsize, text=bkGreen!50!black, anchor=west] at (axis cs:4.45,15) {2\R{回}{かい}};
  \node[font=\sffamily\scriptsize, text=bkOrange!80!black, anchor=west] at (axis cs:4.45,20) {1\R{回}{かい}};
  \node[font=\sffamily\scriptsize, text=bkRed!80!black, anchor=west] at (axis cs:4.45,25) {0\R{回}{かい}};
\end{axis}
\end{tikzpicture}
\end{minipage}\hfill
\begin{minipage}[c]{0.46\linewidth}
\R{平方完成}{へい|ほう|かん|せい}すると、\R{右辺}{う|へん}が「プラス」なら $\pm\sqrt{\phantom{2}}$ で\R{解}{かい}が2つ、「0」なら $t-2=0$ だけで1つ、「マイナス」なら2\R{乗}{じょう}してマイナスになる\R{数}{かず}はないので、\R{解}{かい}はありません（\pageref{bk:3-10:root}ページ）。

グラフでは、\R{横}{よこ}の\R{直線}{ちょく|せん}が\R{放物線}{ほう|ぶつ|せん}と2\R{回}{かい}\R{交}{まじ}わる・1\R{回}{かい}だけふれる・\R{交}{まじ}わらない、に\R{対応}{たい|おう}します。
\end{minipage}

\medskip
\R{解}{かい}の\R{公式}{こう|しき}の $\sqrt{\phantom{2}}$ の\R{中}{なか}にある $b^2-4ac$ も、\R{同}{おな}じ\R{符号}{ふ|ごう}になっています。$b^2-4ac$ を\R{見}{み}るだけで、\R{解}{かい}の\R{個数}{こ|すう}がわかるのです。$b^2-4ac$ を、2\R{次方程式}{じ|ほう|てい|しき}の\textbf{\R{判別式}{はん|べつ|しき}}といい、$D$ で\R{表}{あらわ}します。\bkidx{はんべつしき}{\R{判別式}{はん|べつ|しき}}

\begin{taisetsu}
$ax^2+bx+c=0$ の\R{判別式}{はん|べつ|しき}\quad $D=b^2-4ac$\par\vspace{2pt}
$D>0$ なら\R{解}{かい}は2つ。$D=0$ なら\R{解}{かい}は1つ。$D<0$ なら\R{解}{かい}はない。
\end{taisetsu}
「ボールがその\R{高}{たか}さに\R{届}{とど}くか」は、\R{判別式}{はん|べつ|しき}の\R{符号}{ふ|ごう}で\R{判定}{はん|てい}できます。$D\geqq0$ なら\R{届}{とど}き、$D<0$ なら\R{届}{とど}きません。\R{第}{だい}6\R{巻}{かん}の 6-12 では、この\R{判別式}{はん|べつ|しき}で「ばねにつるしたおもりが、ゆれながら\R{止}{と}まるか、ゆれずに\R{止}{と}まるか」を\R{判定}{はん|てい}します。

\bkthrkhead{\R{解}{と}から、\R{因数分解}{いん|すう|ぶん|かい}の\R{形}{かたち}を\R{作}{つく}る}
$ax^2+bx+c=0$ の\R{解}{かい}が2つあるとき、その\R{解}{かい}を $\alpha$（アルファ）、$\beta$（ベータ）と\R{書}{か}くことがよくあります。\R{解}{かい}が $\alpha$、$\beta$ なら、\R{左辺}{さ|へん}は
\bkthrkline{$ax^2+bx+c=a(x-\alpha)(x-\beta)$}
と\R{因数分解}{いん|すう|ぶん|かい}できます。$x=\alpha$ を\R{入}{い}れると $x-\alpha=0$ で、\R{右辺}{う|へん}が0になるからです。たとえば $2x^2-8x+6=0$ の\R{解}{かい}は1と3（\R{両辺}{りょう|へん}を2でわると $x^2-4x+3=0$）なので、$2x^2-8x+6=2(x-1)(x-3)$ です。グラフ $y=2x^2-8x+6$ は、$x$ \R{軸}{じく}と $x=\alpha$、$x=\beta$ の2\R{点}{てん}で\R{交}{まじ}わります。

\bkkigocards{\bkkigocard{$\alpha$、$\beta$}{アルファ、ベータ}{ギリシャ\R{文字}{も|じ}。2\R{次方程式}{じ|ほう|てい|しき}の2つの\R{解}{かい}によく\R{使}{つか}う。}{alpha, beta}{3-10}}
{\bkkigocard{$D$}{ディー}{\R{判別式}{はん|べつ|しき} $b^2-4ac$。\R{解}{かい}の\R{個数}{こ|すう}がわかる。}{discriminant}{3-10}}

\bksubsection{10}{\R{式}{しき}がプラスになるのは、どこ？}

\leavevmode\label{bk:3-10:fugou}ボールが \qty{15}{m} より\R{高}{たか}い\R{所}{ところ}にいるのは、\R{何秒後}{なん|びょう|ご}から\R{何秒後}{なん|びょう|ご}まででしょう。\R{高}{たか}さ $y$ から15をひいた $y-15$ が、プラスになる\R{時刻}{じ|こく}をさがします。
\[ y-15=-5t^2+20t-15=-5(t^2-4t+3)=-5(t-1)(t-3) \]
$-5$、$t-1$、$t-3$ の3つのかけ\R{算}{ざん}です。それぞれの\R{符号}{ふ|ごう}を\R{表}{ひょう}にすると、かけ\R{算}{ざん}の\R{符号}{ふ|ごう}がわかります（2-2）。5-7 で\R{使}{つか}う\R{表}{ひょう}と\R{同}{おな}じ\R{形}{かたち}です。

\begin{bkcenter}
\begin{tblr}{colspec={Q[c,m,wd=20mm] Q[c,m,wd=16mm] Q[c,m,wd=10mm] Q[c,m,wd=20mm] Q[c,m,wd=10mm] Q[c,m,wd=16mm]}, colsep=4pt, hlines={0.6pt,bkLine}, vline{2}={0.6pt,bkLine}, rowsep=2pt, row{1}={bg=bkTheme!10, font=\small}}
$t$ & $t<1$ & 1 & $1<t<3$ & 3 & $t>3$ \\
$-5$ & $-$ & $-$ & $-$ & $-$ & $-$ \\
$t-1$ & $-$ & 0 & $+$ & $+$ & $+$ \\
$t-3$ & $-$ & $-$ & $-$ & 0 & $+$ \\
$y-15$ & $-$ & 0 & \SetCell{bg=bkYellow!30} $+$ & 0 & $-$ \\
\end{tblr}
\end{bkcenter}
たとえば $1<t<3$ では、$(-)\times(+)\times(-)=(+)$ です。\R{数}{かず}を\R{入}{い}れて\R{確}{たし}かめると、$t=2$ で $-5\times1\times(-1)=5$ で、たしかにプラスです（\R{高}{たか}さ \qty{20}{m} は \qty{15}{m} より \qty{5}{m} \R{高}{たか}い）。

だから、ボールが \qty{15}{m} より\R{高}{たか}いのは、1\R{秒後}{びょう|ご}から3\R{秒後}{びょう|ご}までのあいだです。グラフで、\R{放物線}{ほう|ぶつ|せん}がオレンジの\R{直線}{ちょく|せん}より\R{上}{うえ}にある\R{部分}{ぶ|ぶん}です（\pageref{bk:3-10:toku}ページの\R{図}{ず}）。このように、\R{不等号}{ふ|とう|ごう}のついた2\R{次}{じ}の\R{式}{しき}（$y-15>0$ など）を\R{考}{かんが}えることを、\R{高校}{こう|こう}では「2\R{次不等式}{じ|ふ|とう|しき}」として\R{学}{まな}びます。

\begin{kotoba}
\small
\bkword{2\R{次方程式}{じ|ほう|てい|しき}}{$ax^2+bx+c=0$（$a$ は0でない）の\R{形}{かたち}にまとめられる\R{方程式}{ほう|てい|しき}。}
\bkword{\R{解}{かい}}{\R{方程式}{ほう|てい|しき}をなりたたせる\R{文字}{も|じ}の\R{値}{あたい}。2\R{次方程式}{じ|ほう|てい|しき}の\R{解}{かい}は、2つ・1つ・なしのどれか。}
\bkword{\R{因数分解}{いん|すう|ぶん|かい}・\R{因数}{いん|すう}}{\R{和}{わ}の\R{形}{かたち}の\R{式}{しき}を、かけ\R{算}{ざん}の\R{形}{かたち}に\R{直}{なお}すこと。かけ\R{算}{ざん}を\R{作}{つく}る1つ1つの\R{式}{しき}が\R{因数}{いん|すう}。}
\bkword{\R{解}{かい}の\R{公式}{こう|しき}}{$x=\dfrac{-b\pm\sqrt{b^2-4ac}}{2a}$。\R{平方完成}{へい|ほう|かん|せい}から\R{作}{つく}る。}
\bkword{\R{判別式}{はん|べつ|しき}}{$D=b^2-4ac$。プラスなら\R{解}{かい}2つ、0なら1つ、マイナスならない。}
\bkword{\R{有理化}{ゆう|り|か}}{\R{分母}{ぶん|ぼ}と\R{分子}{ぶん|し}に\R{同}{おな}じ $\sqrt{\phantom{2}}$ をかけて、\R{分母}{ぶん|ぼ}の $\sqrt{\phantom{2}}$ をなくすこと。}
\end{kotoba}

%% ---------------- 例題 ----------------
\Needspace{14\baselineskip}
\bksection{reidai}{\R{例題}{れい|だい}でたしかめよう}

\begin{reidai}[$x^2=k$ の\R{形}{かたち}の\R{方程式}{ほう|てい|しき}]
\R{次}{つぎ}の\R{方程式}{ほう|てい|しき}を\R{解}{かい}きましょう。（3）は\R{時刻}{じ|こく} $t$ \R{秒}{びょう}なので、\R{場面}{ば|めん}に\R{合}{あ}う\R{解}{かい}を\R{答}{こた}えましょう。\par
（1）\nobreak $x^2=25$\quad（2）\nobreak $x^2=7$\quad（3）\nobreak $5t^2=45$\quad（4）\nobreak $(x-2)^2=9$\par
\kangaekata 2\R{乗}{じょう}して $k$ になる\R{数}{かず}は、$\sqrt{k}$ と $-\sqrt{k}$ の2つ。（3）は\R{先}{さき}に\R{両辺}{りょう|へん}を5でわります。（4）は $x-2$ を1つの\R{箱}{はこ}と\R{見}{み}ます。\par
\shiki（1）\nobreak $5^2=25$、$(-5)^2=25$ だから $x=\pm5$\par
\hspace*{2.6\zw}（2）\nobreak 2\R{乗}{じょう}して7になる\R{整数}{せい|すう}はないので $x=\pm\sqrt{7}$\par
\hspace*{2.6\zw}（3）\nobreak \R{両辺}{りょう|へん}を5でわって $t^2=9$。$t=\pm3$。\R{時刻}{じ|こく}は0より\R{後}{あと}なので $t=3$\par
\hspace*{2.6\zw}（4）\nobreak $x-2=\pm3$。$x-2=3$ から $x=5$、$x-2=-3$ から $x=-1$\par
\kotae（1）\nobreak $x=\pm5$\quad（2）\nobreak $x=\pm\sqrt{7}$\quad（3）\nobreak 3\R{秒後}{びょう|ご}\quad（4）\nobreak $x=5,\ -1$
\end{reidai}

\begin{reidai}[かけて0なら、どちらかが0]
\R{次}{つぎ}の\R{方程式}{ほう|てい|しき}を\R{解}{かい}きましょう。\par（1）\nobreak $x(x-2)=0$\quad（2）\nobreak $(t+5)(t-3)=0$\par
\kangaekata かけて0なので、どちらかのかっこ（\R{因数}{いん|すう}）が0です。それぞれを0とおいた1\R{次方程式}{じ|ほう|てい|しき}を\R{解}{かい}きます。\par
\shiki（1）\nobreak $x=0$ または $x-2=0$。$x-2=0$ から $x=2$\par
\hspace*{2.6\zw}（2）\nobreak $t+5=0$ または $t-3=0$。$t=-5$ または $t=3$\par
\hspace*{2.6\zw}\R{確}{たし}かめ：（2）\nobreak で $t=-5$ なら $0\times(-8)=0$、$t=3$ なら $8\times0=0$\par
\kotae（1）\nobreak $x=0,\ 2$\quad（2）\nobreak $t=-5,\ 3$
\end{reidai}

\bkthrkpair{A}{\begin{jibun}
\R{次}{つぎ}の\R{方程式}{ほう|てい|しき}を\R{解}{かい}きましょう。（3）は\R{時刻}{じ|こく} $t$ \R{秒}{びょう}です。\par
（1）\nobreak $x^2=49$\quad（2）\nobreak $x^2=10$\par
（3）\nobreak $5t^2=80$\quad（4）\nobreak $(x+1)^2=16$
\bkjans{（1）\nobreak $x=\pm7$（2）\nobreak $x=\pm\sqrt{10}$（3）\nobreak $t^2=16$、$t=\pm4$。\R{時刻}{じ|こく}なので4\R{秒後}{びょう|ご}（4）\nobreak $x+1=\pm4$。$x=3,\ -5$}
\end{jibun}}{\begin{jibun}
\R{次}{つぎ}の\R{方程式}{ほう|てい|しき}を\R{解}{かい}きましょう。\par
（1）\nobreak $x(x+4)=0$\par
（2）\nobreak $(x-1)(x-6)=0$
\bkjans{（1）\nobreak $x=0$ または $x+4=0$。$x=0,\ -4$（2）\nobreak $x=1,\ 6$}
\end{jibun}}

\begin{reidai}[\R{共通因数}{きょう|つう|いん|すう}でくくって\R{解}{と}く]
\R{次}{つぎ}の\R{方程式}{ほう|てい|しき}を\R{解}{かい}きましょう。\par（1）\nobreak $x^2-2x=0$\quad（2）\nobreak $3t^2-6t=0$\par
\kangaekata \R{両辺}{りょう|へん}を $x$ でわらずに（\pageref{bk:3-10:tsuma}ページ）、\R{共通因数}{きょう|つう|いん|すう}でくくってかけ\R{算}{ざん}の\R{形}{かたち}にします。\par
\shiki（1）\nobreak $x(x-2)=0$。$x=0$ または $x-2=0$\par
\hspace*{2.6\zw}（2）\nobreak $3t(t-2)=0$。3は0でないので、$t=0$ または $t-2=0$\par
\kotae（1）\nobreak $x=0,\ 2$\quad（2）\nobreak $t=0,\ 2$
\end{reidai}

\begin{jibun}
\R{次}{つぎ}の\R{方程式}{ほう|てい|しき}を\R{解}{かい}きましょう。\par
（1）\nobreak $x^2-x=0$\par
（2）\nobreak $5t^2-20t=0$
\bkjans{（1）\nobreak $x(x-1)=0$ で $x=0,\ 1$（2）\nobreak $5t(t-4)=0$ で $t=0,\ 4$}
\end{jibun}

\begin{reidai}[\R{和}{わ}と\R{積}{せき}で\R{因数分解}{いん|すう|ぶん|かい}する]
\R{次}{つぎ}の\R{式}{しき}を\R{因数分解}{いん|すう|ぶん|かい}しましょう。\par
（1）\nobreak $x^2+7x+12$\quad（2）\nobreak $x^2-8x+12$\quad（3）\nobreak $t^2+2t-15$\quad（4）\nobreak $x^2-2x-8$\par
\kangaekata たして $x$ の\R{係数}{けい|すう}、かけて\R{定数項}{てい|すう|こう}になる2つの\R{数}{かず}をさがします。\R{積}{せき}の\R{符号}{ふ|ごう}から、2つの\R{数}{かず}の\R{符号}{ふ|ごう}をしぼります。\par
\shiki（1）\nobreak \R{積}{せき}も\R{和}{わ}もプラス → 2つともプラス。3と4\par
\hspace*{2.6\zw}（2）\nobreak \R{積}{せき}がプラス、\R{和}{わ}がマイナス → 2つともマイナス。$-2$ と $-6$\par
\hspace*{2.6\zw}（3）\nobreak \R{積}{せき}がマイナス → ちがう\R{符号}{ふ|ごう}。\R{和}{わ}がプラス → \R{大}{おお}きいほうがプラス。5と $-3$\par
\hspace*{2.6\zw}（4）\nobreak \R{積}{せき}がマイナス → ちがう\R{符号}{ふ|ごう}。\R{和}{わ}がマイナス → \R{大}{おお}きいほうがマイナス。$-4$ と2\par
\kotae（1）\nobreak $(x+3)(x+4)$\quad（2）\nobreak $(x-2)(x-6)$\quad（3）\nobreak $(t+5)(t-3)$\quad（4）\nobreak $(x-4)(x+2)$\par
（\R{確}{たし}かめ：（4）\nobreak を\R{展開}{てん|かい}すると $x^2+2x-4x-8=x^2-2x-8$。）
\end{reidai}



\begin{jibun}
\R{次}{つぎ}の\R{式}{しき}を\R{因数分解}{いん|すう|ぶん|かい}しましょう。\par
（1）\nobreak $x^2+6x+8$\quad（2）\nobreak $x^2-5x+6$\par
（3）\nobreak $t^2+3t-10$\quad（4）\nobreak $x^2-x-12$
\bkjans{（1）\nobreak $(x+2)(x+4)$（2）\nobreak $(x-2)(x-3)$（3）\nobreak $(t+5)(t-2)$（4）\nobreak $(x-4)(x+3)$}
\end{jibun}

\begin{reidai}[$x^2-a^2$ と、くくってから\R{因数分解}{いん|すう|ぶん|かい}]
\R{次}{つぎ}の\R{式}{しき}を\R{因数分解}{いん|すう|ぶん|かい}しましょう。\par（1）\nobreak $x^2-1$\quad（2）\nobreak $3x^2-3$\quad（3）\nobreak $12x^2-96x+144$\par
\kangaekata まず、\R{共通因数}{きょう|つう|いん|すう}でくくれないかを\R{見}{み}ます。\R{次}{つぎ}に「2\R{乗}{じょう}ひく2\R{乗}{じょう}」か「\R{和}{わ}と\R{積}{せき}」の\R{形}{かたち}かを\R{見}{み}ます。\par
\shiki（1）\nobreak $x^2-1^2$ なので $(x+1)(x-1)$\par
\hspace*{2.6\zw}（2）\nobreak 3でくくって $3(x^2-1)$。かっこの\R{中}{なか}は（1）と\R{同}{おな}じ\par
\hspace*{2.6\zw}（3）\nobreak 12でくくって $12(x^2-8x+12)$。\R{和}{わ}が $-8$、\R{積}{せき}が12 の2つの\R{数}{かず}は $-2$ と $-6$\par
\kotae（1）\nobreak $(x+1)(x-1)$\quad（2）\nobreak $3(x+1)(x-1)$\quad（3）\nobreak $12(x-2)(x-6)$
\end{reidai}

\begin{reidai}[\R{因数分解}{いん|すう|ぶん|かい}で2\R{次方程式}{じ|ほう|てい|しき}を\R{解}{と}く]
\R{次}{つぎ}の\R{方程式}{ほう|てい|しき}を\R{解}{かい}きましょう。（2）（3）は\R{時刻}{じ|こく} $t$ \R{秒}{びょう}（$t>0$）です。\par
（1）\nobreak $x^2-8x+12=0$\quad（2）\nobreak $t^2=15-2t$\quad（3）\nobreak $20t-5t^2=15$\par
\kangaekata \R{移項}{い|こう}して\R{右辺}{う|へん}を0にしてから、\R{左辺}{さ|へん}を\R{因数分解}{いん|すう|ぶん|かい}します。\R{最後}{さい|ご}に、\R{場面}{ば|めん}に\R{合}{あ}う\R{解}{かい}をえらびます。\par
\shiki（1）\nobreak $(x-2)(x-6)=0$ より $x=2,\ 6$\par
\hspace*{2.6\zw}（2）\nobreak $2t$ と $-15$ を\R{左辺}{さ|へん}に\R{移項}{い|こう}して $t^2+2t-15=0$\par
\hspace*{4.3\zw}$(t+5)(t-3)=0$ より $t=-5,\ 3$。$t>0$ なので $t=3$\par
\hspace*{2.6\zw}（3）\nobreak \R{移項}{い|こう}して $-5t^2+20t-15=0$。\R{両辺}{りょう|へん}を $-5$ でわって $t^2-4t+3=0$\par
\hspace*{4.3\zw}$(t-1)(t-3)=0$ より $t=1,\ 3$。どちらも\R{場面}{ば|めん}に\R{合}{あ}う\par
\kotae（1）\nobreak $x=2,\ 6$\quad（2）\nobreak 3\R{秒後}{びょう|ご}\quad（3）\nobreak 1\R{秒後}{びょう|ご}と3\R{秒後}{びょう|ご}
\end{reidai}

\bkthrkpair{C}{\begin{jibun}
\R{次}{つぎ}の\R{式}{しき}を\R{因数分解}{いん|すう|ぶん|かい}しましょう。\par
（1）\nobreak $x^2-9$\quad（2）\nobreak $2x^2-8$\par
（3）\nobreak $3t^2-12t+9$
\bkjans{（1）\nobreak $(x+3)(x-3)$（2）\nobreak $2(x^2-4)=2(x+2)(x-2)$（3）\nobreak $3(t^2-4t+3)=3(t-1)(t-3)$}
\end{jibun}}{\begin{jibun}
\R{次}{つぎ}の\R{方程式}{ほう|てい|しき}を\R{解}{かい}きましょう。（2）（3）は\R{時刻}{じ|こく} $t$ \R{秒}{びょう}（$t>0$）です。\par
（1）\nobreak $x^2+x-6=0$\par
（2）\nobreak $t^2=8-2t$\par
（3）\nobreak $25t-5t^2=20$
\bkjans{（1）\nobreak $(x+3)(x-2)=0$ で $x=-3,\ 2$（2）\nobreak $t^2+2t-8=0$、$(t+4)(t-2)=0$。$t>0$ なので2\R{秒後}{びょう|ご}（3）\nobreak $t^2-5t+4=0$、$(t-1)(t-4)=0$。1\R{秒後}{びょう|ご}と4\R{秒後}{びょう|ご}}
\end{jibun}}

\begin{reidai}[\R{平方完成}{へい|ほう|かん|せい}で\R{解}{と}く]
\R{次}{つぎ}の\R{方程式}{ほう|てい|しき}を\R{解}{かい}きましょう。（2）は\R{時刻}{じ|こく} $t$ \R{秒}{びょう}で、\R{約}{やく}\R{何秒}{なん|びょう}かも\R{答}{こた}えましょう（$\sqrt{2}$ は\R{約}{やく}1.41）。\par
（1）\nobreak $x^2-4x+1=0$\quad（2）\nobreak $20t-5t^2=10$\par
\kangaekata \R{整数}{せい|すう}の\R{和}{わ}と\R{積}{せき}では\R{見}{み}つからないので、\R{平方完成}{へい|ほう|かん|せい}して「（かっこ）の2\R{乗}{じょう}＝\R{数}{かず}」の\R{形}{かたち}にします。\par
\shiki（1）\nobreak $x^2-4x+4-4+1=0$ より $(x-2)^2-3=0$、$(x-2)^2=3$\par
\hspace*{4.3\zw}$x-2=\pm\sqrt{3}$ より $x=2\pm\sqrt{3}$\par
\hspace*{2.6\zw}（2）\nobreak \R{移項}{い|こう}して $-5$ でわると $t^2-4t+2=0$\par
\hspace*{4.3\zw}$(t-2)^2-2=0$ より $(t-2)^2=2$、$t-2=\pm\sqrt{2}$、$t=2\pm\sqrt{2}$\par
\hspace*{4.3\zw}$2-1.41=0.59$、$2+1.41=3.41$\par
\kotae（1）\nobreak $x=2\pm\sqrt{3}$\quad（2）\nobreak $t=2\pm\sqrt{2}$（\R{約}{やく}0.59\R{秒後}{びょう|ご}と\R{約}{やく}3.41\R{秒後}{びょう|ご}）
\end{reidai}

\begin{reidai}[$\sqrt{\phantom{2}}$ の\R{計算}{けい|さん}]
\R{次}{つぎ}の\R{数}{かず}を、（1）（2）は $\sqrt{\phantom{2}}$ の\R{中}{なか}をできるだけ\R{小}{ちい}さい\R{整数}{せい|すう}に、（3）（4）は\R{分母}{ぶん|ぼ}を\R{有理化}{ゆう|り|か}しましょう。\par
（1）\nobreak $\sqrt{12}$\quad（2）\nobreak $\sqrt{18}$\quad（3）\nobreak $\dfrac{1}{\sqrt{2}}$\quad（4）\nobreak $\dfrac{6}{\sqrt{3}}$\par
\kangaekata （1）（2）は「2\R{乗}{じょう}の\R{数}{かず}×ほかの\R{数}{かず}」に\R{分}{わ}けます。（3）（4）は\R{分母}{ぶん|ぼ}と\R{分子}{ぶん|し}に\R{同}{おな}じ $\sqrt{\phantom{2}}$ をかけます。\par
\shiki（1）\nobreak $\sqrt{4\times3}=\sqrt{4}\times\sqrt{3}=2\sqrt{3}$\quad（2）\nobreak $\sqrt{9\times2}=3\sqrt{2}$\par
\hspace*{2.6\zw}（3）\nobreak $\dfrac{1\times\sqrt{2}}{\sqrt{2}\times\sqrt{2}}=\dfrac{\sqrt{2}}{2}$\quad（4）\nobreak $\dfrac{6\times\sqrt{3}}{\sqrt{3}\times\sqrt{3}}=\dfrac{6\sqrt{3}}{3}=2\sqrt{3}$\par
\kotae（1）\nobreak $2\sqrt{3}$\quad（2）\nobreak $3\sqrt{2}$\quad（3）\nobreak $\dfrac{\sqrt{2}}{2}$\quad（4）\nobreak $2\sqrt{3}$\par
（4）は、\R{最後}{さい|ご}に\R{分母}{ぶん|ぼ}の3と\R{分子}{ぶん|し}の6を\R{約分}{やく|ぶん}しました。\R{電卓}{でん|たく}で\R{確}{たし}かめると、\par
$6\div1.732\cdots=3.464\cdots$、$2\times1.732\cdots=3.464\cdots$ で\R{合}{あ}います。
\end{reidai}

\bkthrkpair{D}{\begin{jibun}
\R{次}{つぎ}の\R{方程式}{ほう|てい|しき}を\R{解}{かい}きましょう。（2）は\R{時刻}{じ|こく} $t$ \R{秒}{びょう}で、\R{約}{やく}\R{何秒}{なん|びょう}かも\R{答}{こた}えましょう（$\sqrt{3}$ は\R{約}{やく}1.73）。\par
（1）\nobreak $x^2+6x+4=0$\par
（2）\nobreak $20t-5t^2=5$
\bkjans{（1）\nobreak $(x+3)^2=5$、$x=-3\pm\sqrt{5}$（2）\nobreak $t^2-4t+1=0$、$(t-2)^2=3$、$t=2\pm\sqrt{3}$。\R{約}{やく}0.27\R{秒後}{びょう|ご}と\R{約}{やく}3.73\R{秒後}{びょう|ご}}
\end{jibun}}{\begin{jibun}
（1）（2）は $\sqrt{\phantom{2}}$ の\R{中}{なか}を\R{小}{ちい}さく、（3）（4）は\R{分母}{ぶん|ぼ}を\R{有理化}{ゆう|り|か}しましょう。\par
（1）\nobreak $\sqrt{8}$\quad（2）\nobreak $\sqrt{27}$\par
（3）\nobreak $\dfrac{1}{\sqrt{3}}$\quad（4）\nobreak $\dfrac{4}{\sqrt{2}}$
\bkjans{（1）\nobreak $2\sqrt{2}$（2）\nobreak $3\sqrt{3}$（3）\nobreak $\dfrac{\sqrt{3}}{3}$（4）\nobreak $\dfrac{4\sqrt{2}}{2}=2\sqrt{2}$}
\end{jibun}}

\begin{reidai}[\R{解}{かい}の\R{公式}{こう|しき}で\R{解}{と}く]
\R{次}{つぎ}の\R{方程式}{ほう|てい|しき}を、\R{解}{かい}の\R{公式}{こう|しき}で\R{解}{かい}きましょう。\par（1）\nobreak $2x^2+3x-1=0$\quad（2）\nobreak $x^2-6x+4=0$\par
\kangaekata $a$、$b$、$c$ を\R{符号}{ふ|ごう}ごと\R{読}{よ}みとって、\R{公式}{こう|しき}に\R{入}{い}れます。$\sqrt{\phantom{2}}$ の\R{中}{なか}が\R{小}{ちい}さくできるときは\R{小}{ちい}さくし、\R{約分}{やく|ぶん}できるときは\R{約分}{やく|ぶん}します。\par
\shiki（1）\nobreak $a=2$、$b=3$、$c=-1$\par
\addvspace{3pt}\hspace*{4.3\zw}$x=\dfrac{-3\pm\sqrt{3^2-4\times2\times(-1)}}{2\times2}=\dfrac{-3\pm\sqrt{9+8}}{4}=\dfrac{-3\pm\sqrt{17}}{4}$\par\vspace{2pt}
\hspace*{2.6\zw}（2）\nobreak $a=1$、$b=-6$、$c=4$\par
\addvspace{3pt}\hspace*{4.3\zw}$x=\dfrac{6\pm\sqrt{(-6)^2-4\times1\times4}}{2\times1}=\dfrac{6\pm\sqrt{36-16}}{2}=\dfrac{6\pm\sqrt{20}}{2}$\par\addvspace{3pt}
\hspace*{5.3\zw}$=\dfrac{6\pm2\sqrt{5}}{2}=3\pm\sqrt{5}$\par\vspace{2pt}
\hspace*{4.3\zw}（$\sqrt{20}=2\sqrt{5}$。\R{分子}{ぶん|し}の6も $2\sqrt{5}$ も2でわれる）\par
\kotae（1）\nobreak $x=\dfrac{-3\pm\sqrt{17}}{4}$\quad（2）\nobreak $x=3\pm\sqrt{5}$
\end{reidai}

\begin{jibun}
\R{次}{つぎ}の\R{方程式}{ほう|てい|しき}を、\R{解}{かい}の\R{公式}{こう|しき}で\R{解}{かい}きましょう。\par
（1）\nobreak $3x^2+5x+1=0$\par
（2）\nobreak $x^2-2x-2=0$
\bkjans{（1）\nobreak $x=\dfrac{-5\pm\sqrt{25-12}}{6}=\dfrac{-5\pm\sqrt{13}}{6}$（2）\nobreak $x=\dfrac{2\pm\sqrt{4+8}}{2}=\dfrac{2\pm2\sqrt{3}}{2}=1\pm\sqrt{3}$}
\end{jibun}

\begin{reidai}[\R{判別式}{はん|べつ|しき}で、\R{解}{かい}の\R{個数}{こ|すう}を\R{調}{しら}べる]
\R{次}{つぎ}の2\R{次方程式}{じ|ほう|てい|しき}の\R{解}{かい}の\R{個数}{こ|すう}を、\R{判別式}{はん|べつ|しき}で\R{調}{しら}べましょう。（4）は\R{答}{こた}えも\R{求}{もと}めましょう。\par
（1）\nobreak $x^2-4x+3=0$\quad（2）\nobreak $x^2-4x+4=0$\quad（3）\nobreak $x^2-4x+5=0$\par
（4）\nobreak \qty{10}{m/s} で\R{真上}{ま|うえ}に\R{投}{な}げたボールの\R{高}{たか}さは $y=10t-5t^2$（m）です。\qty{6}{m} の\R{高}{たか}さに\R{届}{とど}きますか。\par
\kangaekata $D=b^2-4ac$ の\R{符号}{ふ|ごう}を\R{調}{しら}べます。（4）は $10t-5t^2=6$ を $ax^2+bx+c=0$ の\R{形}{かたち}にします。\par
\shiki（1）\nobreak $D=(-4)^2-4\times1\times3=16-12=4>0$\par
\hspace*{2.6\zw}（2）\nobreak $D=(-4)^2-4\times1\times4=16-16=0$\par
\hspace*{2.6\zw}（3）\nobreak $D=(-4)^2-4\times1\times5=16-20=-4<0$\par
\hspace*{2.6\zw}（4）\nobreak \R{移項}{い|こう}して $-5t^2+10t-6=0$。\par
\hspace*{4.3\zw}\R{両辺}{りょう|へん}に $-1$ をかけて $5t^2-10t+6=0$\par
\hspace*{4.3\zw}$D=(-10)^2-4\times5\times6=100-120=-20<0$\par
\kotae（1）\nobreak 2つ\quad（2）\nobreak 1つ\quad（3）\nobreak ない\quad（4）\nobreak \R{解}{かい}がないので\R{届}{とど}かない\par
（（4）：\R{平方完成}{へい|ほう|かん|せい}すると $y=-5(t-1)^2+5$ で、\R{最高点}{さい|こう|てん}は \qty{5}{m}。\qty{6}{m} には\R{届}{とど}かないことと\R{合}{あ}う。）
\end{reidai}



\begin{jibun}
\R{解}{かい}の\R{個数}{こ|すう}を\R{判別式}{はん|べつ|しき}で\R{調}{しら}べましょう。\par
（1）\nobreak $2x^2+3x+2=0$\quad（2）\nobreak $x^2+6x+9=0$\par
（3）\nobreak $x^2+x-1=0$\par
（4）\nobreak \qty{20}{m/s} で\R{真上}{ま|うえ}に\R{投}{な}げたボール（$y=20t-5t^2$）は、\qty{18}{m} の\R{高}{たか}さに\R{届}{とど}きますか。
\bkjans{（1）\nobreak $D=9-16=-7$ でない（2）\nobreak $D=36-36=0$ で1つ（3）\nobreak $D=1+4=5$ で2つ（4）\nobreak $5t^2-20t+18=0$、$D=400-360=40>0$ なので\R{届}{とど}く（\R{上}{のぼ}りと\R{下}{くだ}りの2\R{回}{かい}）}
\end{jibun}

\begin{reidai}[\R{式}{しき}がプラスになるところ]
（1）\nobreak $(x-1)(x-3)$ がプラスになる $x$ と、マイナスになる $x$ の\R{範囲}{はん|い}を\R{答}{こた}えましょう。\par
（2）\nobreak \qty{20}{m/s} で\R{真上}{ま|うえ}に\R{投}{な}げたボール（$y=20t-5t^2$）が、\qty{15}{m} より\R{高}{たか}い\R{所}{ところ}にいるのは、\R{何秒後}{なん|びょう|ご}から\R{何秒後}{なん|びょう|ご}までですか。\par
\kangaekata \R{因数}{いん|すう}ごとの\R{符号}{ふ|ごう}を\R{表}{ひょう}にして、かけ\R{算}{ざん}の\R{符号}{ふ|ごう}を\R{調}{しら}べます。（2）は $y-15$ を\R{因数分解}{いん|すう|ぶん|かい}します。\par
\shiki（1）\nobreak $x<1$ では $(-)\times(-)=(+)$、$1<x<3$ では $(+)\times(-)=(-)$、$x>3$ では $(+)\times(+)=(+)$\par
\hspace*{2.6\zw}（2）\nobreak $y-15=-5t^2+20t-15=-5(t-1)(t-3)$\par
\hspace*{4.3\zw}$1<t<3$ では $(-)\times(+)\times(-)=(+)$ なので、$y-15$ がプラス\par
\kotae（1）\nobreak プラスは $x<1$ と $x>3$、マイナスは $1<x<3$\quad（2）\nobreak 1\R{秒後}{びょう|ご}から3\R{秒後}{びょう|ご}まで
\end{reidai}

\begin{jibun}
（1）\nobreak $(x+2)(x-4)$ がプラスになる $x$ と、マイナスになる $x$ の\R{範囲}{はん|い}を\R{答}{こた}えましょう。\quad（2）\nobreak \qty{25}{m/s} で\R{真上}{ま|うえ}に\R{投}{な}げたボール（$y=25t-5t^2$）が、\qty{20}{m} より\R{高}{たか}い\R{所}{ところ}にいるのは、\R{何秒後}{なん|びょう|ご}から\R{何秒後}{なん|びょう|ご}までですか。
\bkjans{（1）\nobreak プラスは $x<-2$ と $x>4$、マイナスは $-2<x<4$（2）\nobreak $y-20=-5(t-1)(t-4)$ で、1\R{秒後}{びょう|ご}から4\R{秒後}{びょう|ご}まで}
\end{jibun}

%% ---------------- まちがえやすいところ・物理工学 ----------------
\Needspace{14\baselineskip}
\bksection{tsumazuki}{まちがえやすいところをチェック}\label{bk:3-10:tsuma}

\begin{tsumazuki}{\R{両辺}{りょう|へん}を $x$ でわってしまう}
$x^2=2x$ を\R{解}{と}くと？
\bkng{\R{両辺}{りょう|へん}を $x$ でわって、$x=2$}
\bkok{$x^2-2x=0$、$x(x-2)=0$ で、$x=0,\ 2$}
$x=0$ のとき、0でわることになるので、わってはいけません。$x=0$ を\R{入}{い}れると $0=0$ でなりたつのに、わると、この\R{解}{かい}が\R{消}{け}えてしまいます。
\end{tsumazuki}

\begin{tsumazuki}{$-3^2$ と $(-3)^2$}
$x=-3$ が $x^2+2x-3=0$ の\R{解}{と}か、\R{確}{たし}かめると？
\bkng{$-3^2+2\times(-3)-3=-9-6-3=-18$ で、\R{解}{かい}ではない}
\bkok{$(-3)^2+2\times(-3)-3=9-6-3=0$ で、\R{解}{かい}である}
$x$ にマイナスの\R{数}{かず}を\R{入}{い}れるときは、かっこごと\R{入}{い}れます（2-13）。$-3^2$ は「3の2\R{乗}{じょう}にマイナスをつけた\R{数}{かず}」で $-9$ です。
\end{tsumazuki}

\begin{tsumazuki}{かけて0の\R{解}{かい}の\R{符号}{ふ|ごう}}
\bkng{$(x-2)(x-6)=0$ だから、$x=-2,\ -6$}
\bkok{$x-2=0$ から $x=2$、$x-6=0$ から $x=6$}
「かっこの\R{中}{なか}が0になる $x$」が\R{解}{かい}です。\R{答}{こた}えを\R{入}{い}れて、かっこの\R{中}{なか}が0になるか\R{確}{たし}かめましょう。
\end{tsumazuki}

\begin{tsumazuki}{0でないのに「どちらか」を\R{使}{つか}う}
\bkng{$(x-1)(x-3)=3$ だから、$x-1=3$ または $x-3=3$}
\bkok{\R{展開}{てん|かい}して $x^2-4x+3=3$、$x^2-4x=0$、$x(x-4)=0$ で $x=0,\ 4$}
「どちらかが\R{決}{き}まる」のは、かけて0のときだけです。\R{先}{さき}に\R{右辺}{う|へん}を0にしましょう。
\end{tsumazuki}

\begin{tsumazuki}{\R{解}{かい}を1つだけ\R{答}{こた}える}
\bkng{$x^2=9$ だから、$x=3$}
\bkok{$x=\pm3$（$3$ と $-3$）}
$(-3)^2=9$ も\R{忘}{わす}れずに。\R{場面}{ば|めん}に\R{合}{あ}わない\R{解}{かい}をすてるのは、\R{両方}{りょう|ほう}を\R{出}{だ}してからです。
\end{tsumazuki}

\begin{tsumazuki}{\R{解}{かい}の\R{公式}{こう|しき}で、$b$ の\R{符号}{ふ|ごう}をまちがえる}
$x^2-4x+1=0$ を\R{解}{かい}の\R{公式}{こう|しき}で\R{解}{と}くと？
\bkng{$x=\dfrac{-4\pm\sqrt{-16-4}}{2}$}
\bkok{$x=\dfrac{4\pm\sqrt{16-4}}{2}=\dfrac{4\pm2\sqrt{3}}{2}=2\pm\sqrt{3}$}
$b=-4$ なので、$-b=-(-4)=4$、$b^2=(-4)^2=16$ です。$b$ は\R{符号}{ふ|ごう}ごと\R{読}{よ}みとり、かっこをつけて\R{入}{い}れましょう。
\end{tsumazuki}

\begin{butsuri}
\par\addvspace{3pt}\textbf{\R{家}{いえ}のコンセントの $\sqrt{2}$}\par\noindent
\noindent\begin{minipage}[c]{0.58\linewidth}
\R{家}{いえ}のコンセントの\R{電気}{でん|き}（\R{交流}{こう|りゅう}）は、\R{電圧}{でん|あつ}が1\R{秒}{びょう}に50\R{回}{かい}または60\R{回}{かい}、プラスとマイナスに\R{入}{い}れかわっています。「\qty{100}{V}」というのは、ずっと\R{同}{おな}じ\R{向}{む}きに\R{流}{なが}れる\R{電気}{でん|き}（\R{直流}{ちょく|りゅう}）の \qty{100}{V} と\R{同}{おな}じ\R{働}{はたら}きをする、という\R{意味}{い|み}の\R{値}{あたい}（\R{実効値}{じっ|こう|ち}）です。いちばん\R{高}{たか}い\R{所}{ところ}の\R{電圧}{でん|あつ}は、その $\sqrt{2}$ \R{倍}{ばい}の $100\sqrt{2}$、\R{約}{やく} \qty{141}{V} あります。\R{反対}{はん|たい}に、\R{実効値}{じっ|こう|ち}は\R{最大値}{さい|だい|ち}に $\dfrac{1}{\sqrt{2}}=\dfrac{\sqrt{2}}{2}$（\R{約}{やく}0.707）をかけた\R{値}{あたい}です。\R{電気}{でん|き}の\R{設計}{せっ|けい}では、\R{有理化}{ゆう|り|か}した $\dfrac{\sqrt{2}}{2}$ が\R{毎日}{まい|にち}のように\R{出}{で}てきます。
\end{minipage}\hfill
\begin{minipage}[c]{0.38\linewidth}\centering
\begin{tikzpicture}
\begin{axis}[bkplot, width=48mm, height=40mm, xmin=0, xmax=2.3, ymin=-170, ymax=190, xtick=\empty, ytick={-141,141},
  ylabel={V}, xlabel={$t$}]
  \addplot[bkcurve, domain=0:2, samples=80] {141*sin(360*x)};
  \addplot[bkcurve2, line width=1pt, dash pattern=on 3pt off 2pt, domain=0:2] {100};
  \node[font=\sffamily\scriptsize, text=bkOrange!80!black, anchor=west] at (axis cs:2.02,100) {100};
\end{axis}
\end{tikzpicture}
\end{minipage}
\end{butsuri}

\begin{butsuri}
\par\addvspace{3pt}\textbf{\R{車}{くるま}が\R{止}{と}まれる\R{速}{はや}さ}\label{bk:3-10:brake}\par\noindent
\noindent\begin{minipage}[c]{0.58\linewidth}
\R{運転}{うん|てん}する\R{人}{ひと}が\R{危}{あぶ}ないと\R{気}{き}づいてからブレーキをふむまで、\R{約}{やく}1\R{秒}{びょう}かかります（\R{反応時間}{はん|のう|じ|かん}）。そのあいだ、\R{車}{くるま}は\R{同}{おな}じ\R{速}{はや}さで\R{進}{すす}みます（\R{空走距離}{くう|そう|きょ|り}）。ブレーキがききはじめると、ぬれた\R{道路}{どう|ろ}では1\R{秒}{びょう}ごとに\R{約}{やく} \qty{5}{m/s} ずつおそくなり、3-9 で\R{見}{み}たように、\R{速}{はや}さの2\R{乗}{じょう}に\R{比例}{ひ|れい}する\R{距離}{きょ|り}で\R{止}{と}まります（\R{制動距離}{せい|どう|きょ|り}）。\R{速}{はや}さを $v$〔m/s〕とすると、\R{止}{と}まるまでの\R{距離}{きょ|り} $d$〔m〕は
\[ d=v\times1+\frac{v^2}{10} \]
です（$v^2$ を10でわるのは、\R{三角形}{さん|かく|けい}の\R{面積}{めん|せき} $v\times\dfrac{v}{5}\div2$ だから。$\dfrac{v}{5}$ はブレーキをかけてから\R{止}{と}まるまでの\R{時間}{じ|かん}）。
\end{minipage}\hfill
\begin{minipage}[c]{0.38\linewidth}\centering
\begin{tikzpicture}[x=1mm, y=1mm]
  \bkcar{(0,4)}{0.42}{bkBlue}
  \fill[bkOrange!35] (0,0) rectangle (20,2.6);
  \fill[bkRed!35] (20,0) rectangle (40,2.6);
  \draw[bkInk, line width=0.5pt] (0,0) rectangle (40,2.6) (20,0) -- (20,2.6);
  \node[font=\sffamily\scriptsize, anchor=north, align=center] at (10,-1) {\R{空走}{くう|そう}\\\qty{10}{m}};
  \node[font=\sffamily\scriptsize, anchor=north, align=center] at (30,-1) {\R{制動}{せい|どう}\\\qty{10}{m}};
  \node[font=\sffamily\scriptsize, anchor=north] at (20,-9.5) {\qty{10}{m/s} のとき \qty{20}{m}};
\end{tikzpicture}
\end{minipage}

\vspace{2pt}
「\qty{20}{m} \R{先}{さき}に\R{人}{ひと}がいるとき、\R{止}{と}まれるのは\R{何}{なん} m/s までか」を\R{考}{かんが}えます。$d=20$ とおくと、
\[ v+\frac{v^2}{10}=20 \]
という2\R{次方程式}{じ|ほう|てい|しき}です。\R{両辺}{りょう|へん}に10をかけて $10v+v^2=200$、\R{移項}{い|こう}して $v^2+10v-200=0$。\R{和}{わ}が10、\R{積}{せき}が $-200$ の2つの\R{数}{かず}は20と $-10$ なので、$(v+20)(v-10)=0$。\R{速}{はや}さはプラスなので $v=10$ です。
\bkthrkline{$\qty{10}{m/s}\times\qty{3600}{s/h}\div\qty{1000}{m/km}=\qty{36}{km/h}$}
\R{時速}{じ|そく} \qty{36}{km} までなら、\qty{20}{m} で\R{止}{と}まれます。
\end{butsuri}

\begin{butsuri}
\par\addvspace{3pt}\textbf{ゆれるか、ゆれないか：\R{特性方程式}{とく|せい|ほう|てい|しき}（6-12 へ）}\par\noindent
\noindent\begin{minipage}[c]{0.6\linewidth}
\R{車}{くるま}のタイヤの\R{近}{ちか}くには、ばねと、ゆれを\R{止}{と}める\R{部品}{ぶ|ひん}（ダンパー）がついています。でこぼこを\R{通}{とお}ったあと、\R{車体}{しゃ|たい}がいつまでもゆれていたら、\R{乗}{じょう}り\R{心地}{ごこ|ち}がよくありません。

\R{第}{だい}6\R{巻}{かん}では、\R{車体}{しゃ|たい}の\R{動}{うご}きを\R{表}{あらわ}す\R{式}{しき}から、たとえば $s^2+3s+2=0$ や $s^2+2s+5=0$ のような2\R{次方程式}{じ|ほう|てい|しき}（\textbf{\R{特性方程式}{とく|せい|ほう|てい|しき}}）が\R{出}{で}てきます。1つめは $(s+1)(s+2)=0$ で\R{解}{かい}が2つあり（$D=9-8=1>0$）、\R{車体}{しゃ|たい}はゆれずに、ゆっくりもとの\R{高}{たか}さにもどります。2つめは $D=4-20=-16<0$ で、この\R{巻}{かん}の\R{数}{かず}の\R{中}{なか}に\R{解}{かい}がなく、\R{車体}{しゃ|たい}はゆれながらもどります。\R{判別式}{はん|べつ|しき}の\R{符号}{ふ|ごう}が、\R{乗}{じょう}り\R{心地}{ごこ|ち}を\R{決}{き}めているのです。ゆれずに、いちばん\R{早}{はや}くもとにもどるのは、ちょうど $D=0$ のときです。\R{実際}{じっ|さい}の\R{車}{くるま}は、\R{乗}{じょう}り\R{心地}{ごこ|ち}のために、ほんの\R{少}{すこ}しだけゆれる（$D$ が\R{少}{すこ}しマイナスの）ように\R{設計}{せっ|けい}することが\R{多}{おお}いのです。
\end{minipage}\hfill
\begin{minipage}[c]{0.36\linewidth}\centering
\begin{tikzpicture}
\begin{axis}[bkplot, width=48mm, height=62mm, xmin=0, xmax=6.6, ymin=-0.55, ymax=1.15, xtick=\empty, ytick=\empty, xlabel={$t$}, ylabel={\R{高}{たか}さ}]
  \addplot[bkcurve, domain=0:6.2, samples=80] {2*exp(-x)-exp(-2*x)};
  \addplot[bkRed, line width=1.4pt, domain=0:6.2, samples=120] {exp(-x)*(cos(deg(2*x))+0.5*sin(deg(2*x)))};
  \node[font=\sffamily\scriptsize, text=bkBlue!80!black, anchor=south west] at (axis cs:1.6,0.62) {$D>0$};
  \node[font=\sffamily\scriptsize, text=bkRed!80!black, anchor=north west] at (axis cs:1.75,-0.25) {$D<0$};
\end{axis}
\end{tikzpicture}
\end{minipage}
\end{butsuri}

%% ---------------- 章末問題 ----------------
\begin{yattemiyou}[\R{表計算}{ひょう|けい|さん}で、2\R{次方程式}{じ|ほう|てい|しき}の\R{解}{かい}をさがそう]
\pageref{bk:3-10:heihou}ページで、\R{高}{たか}さ \qty{10}{m} を\R{通}{とお}る\R{時刻}{じ|こく}は $t=2\pm\sqrt{2}$（\R{約}{やく}0.59\R{秒}{びょう}と\R{約}{やく}3.41\R{秒}{びょう}）と\R{求}{もと}めました。\R{表計算}{ひょう|けい|さん}（2-18）で、この\R{答}{こた}えを\R{確}{たし}かめてみましょう。
\begin{bkcenter}
\bksheet{2}{30}{%
1 & \R{時刻}{じ|こく} $t$ & \R{高}{たか}さ \\
2 & 0 & \texttt{=20*A2-5*A2\^{}2} \\
3 & \texttt{=A2+0.1} & （\R{上}{うえ}をコピー） \\}
\end{bkcenter}
\begin{enumerate}[label={\bkcircled[bkGreen]{\arabic*}}, leftmargin=2.4\zw, itemsep=1pt, topsep=3pt]
\item A\R{列}{れつ}に0から4まで0.1ずつふえる\R{時刻}{じ|こく}を\R{作}{つく}り、B\R{列}{れつ}に\R{高}{たか}さの\R{式}{しき}を\R{下}{した}へコピーしよう。
\item \R{高}{たか}さが10をまたぐ\R{所}{ところ}をさがそう。0.5\R{秒}{びょう}では8.75、0.6\R{秒}{びょう}では10.2 なので、\R{解}{かい}は0.5と0.6のあいだにある。
\item その\R{間}{あいだ}を、こんどは0.01ずつにして\R{調}{しら}べよう。0.58\R{秒}{びょう}で9.918、0.59\R{秒}{びょう}で10.0595 になる。$2-\sqrt{2}=0.5857\cdots$ は、たしかにそのあいだにあるかな？　3.41\R{秒}{びょう}のほうも\R{調}{しら}べよう。
\end{enumerate}
\R{表計算}{ひょう|けい|さん}では、\R{間}{あいだ}をどんどん\R{細}{ほそ}かくすれば、\R{解}{かい}にいくらでも\R{近}{ちか}づけます。\R{答}{こた}えがきれいな\R{形}{かたち}にならない\R{方程式}{ほう|てい|しき}でも、この\R{方法}{ほう|ほう}なら\R{解}{かい}の\R{近}{ちか}くの\R{値}{あたい}がわかります（コンピュータで\R{方程式}{ほう|てい|しき}を\R{解}{と}くときの\R{考}{かんが}え\R{方}{かた}の1つです）。
\end{yattemiyou}

\bkendproblems
\bkproblemgroup{\R{基本}{き|ほん}}{2}

\bkproblem \R{次}{つぎ}の\R{方程式}{ほう|てい|しき}を\R{解}{かい}きましょう。\par\hspace*{1\zw}（1）\nobreak $x^2=36$\quad（2）\nobreak $x^2=5$\quad（3）\nobreak $(x-3)^2=4$
\bkpans{（1）\nobreak $x=\pm6$（2）\nobreak $x=\pm\sqrt{5}$（3）\nobreak $x-3=\pm2$ で $x=5,\ 1$}

\bkproblem \R{次}{つぎ}の\R{式}{しき}を\R{因数分解}{いん|すう|ぶん|かい}しましょう。\par\hspace*{1\zw}（1）\nobreak $x^2+5x+6$\quad（2）\nobreak $x^2-7x+10$\quad（3）\nobreak $x^2+x-12$\par\hspace*{1\zw}（4）\nobreak $x^2-16$\quad（5）\nobreak $2x^2-18$\quad（6）\nobreak $x^2-6x$
\bkpans{（1）\nobreak $(x+2)(x+3)$（2）\nobreak $(x-2)(x-5)$（3）\nobreak $(x+4)(x-3)$（4）\nobreak $(x+4)(x-4)$（5）\nobreak $2(x+3)(x-3)$（6）\nobreak $x(x-6)$}

\bkproblem \R{次}{つぎ}の\R{方程式}{ほう|てい|しき}を\R{解}{かい}きましょう。\par\hspace*{1\zw}（1）\nobreak $x^2-3x=0$\quad（2）\nobreak $x^2-6x+8=0$\quad（3）\nobreak $t^2+t-6=0$
\bkpans{（1）\nobreak $x(x-3)=0$ で $x=0,\ 3$（2）\nobreak $(x-2)(x-4)=0$ で $x=2,\ 4$（3）\nobreak $(t+3)(t-2)=0$ で $t=-3,\ 2$}

\bkproblem \R{次}{つぎ}の\R{方程式}{ほう|てい|しき}を\R{解}{かい}きましょう。（1）は\R{平方完成}{へい|ほう|かん|せい}で、（2）は\R{解}{かい}の\R{公式}{こう|しき}で。\par\hspace*{1\zw}（1）\nobreak $x^2-2x-1=0$\quad（2）\nobreak $2x^2+x-2=0$
\bkpans{（1）\nobreak $(x-1)^2=2$ で $x=1\pm\sqrt{2}$（2）\nobreak $x=\dfrac{-1\pm\sqrt{1+16}}{4}=\dfrac{-1\pm\sqrt{17}}{4}$}

\bkproblem （1）（2）は $\sqrt{\phantom{2}}$ の\R{中}{なか}を\R{小}{ちい}さく、（3）は\R{分母}{ぶん|ぼ}を\R{有理化}{ゆう|り|か}しましょう。\par\hspace*{1\zw}（1）\nobreak $\sqrt{20}$\quad（2）\nobreak $\sqrt{45}$\quad（3）\nobreak $\dfrac{3}{\sqrt{3}}$
\bkpans{（1）\nobreak $2\sqrt{5}$（2）\nobreak $3\sqrt{5}$（3）\nobreak $\dfrac{3\sqrt{3}}{3}=\sqrt{3}$}

\bkproblem \R{次}{つぎ}の2\R{次方程式}{じ|ほう|てい|しき}の\R{解}{かい}の\R{個数}{こ|すう}を、\R{判別式}{はん|べつ|しき}で\R{調}{しら}べましょう。\par\hspace*{1\zw}（1）\nobreak $x^2+2x+3=0$\quad（2）\nobreak $4x^2-4x+1=0$\quad（3）\nobreak $x^2-3x-1=0$
\bkpans{（1）\nobreak $D=4-12=-8$ でない（2）\nobreak $D=16-16=0$ で1つ（3）\nobreak $D=9+4=13$ で2つ}

\bkproblemgroup{チャレンジ}{3}

\bkproblem \qty{30}{m/s} で\R{真上}{ま|うえ}に\R{投}{な}げたボールの、$t$ \R{秒後}{びょう|ご}の\R{高}{たか}さは $y=30t-5t^2$（m）です。（1）\nobreak \R{高}{たか}さ \qty{25}{m} を\R{通}{とお}るのは\R{何秒後}{なん|びょう|ご}ですか。（2）\nobreak \qty{25}{m} より\R{高}{たか}い\R{所}{ところ}にいるのは、\R{何秒後}{なん|びょう|ご}から\R{何秒後}{なん|びょう|ご}までですか。（3）\nobreak \R{地面}{じ|めん}にもどるのは\R{何秒後}{なん|びょう|ご}ですか。
\bkpans{（1）\nobreak $t^2-6t+5=0$、$(t-1)(t-5)=0$ で1\R{秒後}{びょう|ご}（\R{上}{のぼ}り）と5\R{秒後}{びょう|ご}（\R{下}{くだ}り）（2）\nobreak $y-25=-5(t-1)(t-5)$ で、1\R{秒後}{びょう|ご}から5\R{秒後}{びょう|ご}まで（3）\nobreak $-5t^2+30t=0$、$-5t(t-6)=0$ で $t=0,\ 6$。6\R{秒後}{びょう|ご}}

\bkproblem \R{止}{と}まるまでの\R{距離}{きょ|り}が $d=v+\dfrac{v^2}{10}$（m）の\R{車}{くるま}があります（$v$ は\R{速}{はや}さ〔m/s〕）。\qty{60}{m} で\R{止}{と}まれる\R{速}{はや}さは\R{何}{なん} m/s ですか。\R{時速}{じ|そく}\R{何}{なん} km ですか（\pageref{bk:3-10:brake}ページの\R{物理}{ぶつ|り}・\R{工学}{こう|がく}ではこう\R{使}{つか}うと\R{同}{おな}じ\R{手順}{て|じゅん}で、\R{両辺}{りょう|へん}に10をかけてから\R{移項}{い|こう}しよう）。
\bkpans{$v^2+10v-600=0$、$(v+30)(v-20)=0$。$v>0$ なので \qty{20}{m/s}。$\qty{20}{m/s}\times\qty{3600}{s/h}\div\qty{1000}{m/km}=\qty{72}{km/h}$}

\bkproblem \qty{20}{m/s} で\R{真上}{ま|うえ}に\R{投}{な}げたボール（$y=20t-5t^2$）は、\qty{21}{m} の\R{高}{たか}さに\R{届}{とど}きますか。\R{判別式}{はん|べつ|しき}を\R{使}{つか}って\R{答}{こた}え、そのわけを3-9 の\R{最高点}{さい|こう|てん}でも\R{確}{たし}かめましょう。
\bkpans{$5t^2-20t+21=0$、$D=400-420=-20<0$ で\R{届}{とど}かない。\R{最高点}{さい|こう|てん}は \qty{20}{m}（2\R{秒後}{びょう|ご}）なので、\qty{21}{m} には\R{届}{とど}かないことと\R{合}{あ}う}

\bkproblem \qty{30}{m/s} で\R{真上}{ま|うえ}に\R{投}{な}げたボール（$y=30t-5t^2$）が、\R{高}{たか}さ \qty{30}{m} を\R{通}{とお}るのは\R{何秒後}{なん|びょう|ご}ですか。\R{平方完成}{へい|ほう|かん|せい}で\R{解}{かい}き、\R{約}{やく}\R{何秒}{なん|びょう}かも\R{答}{こた}えましょう（$\sqrt{3}$ は\R{約}{やく}1.73）。
\bkpans{$t^2-6t+6=0$、$(t-3)^2=3$、$t=3\pm\sqrt{3}$。\R{約}{やく}1.27\R{秒後}{びょう|ご}（\R{上}{のぼ}り）と\R{約}{やく}4.73\R{秒後}{びょう|ご}（\R{下}{くだ}り）}

\begin{furikaeri}
\begin{checklist}
\item $x^2=k$ の\R{形}{かたち}や\R{因数分解}{いん|すう|ぶん|かい}で、2\R{次方程式}{じ|ほう|てい|しき}を\R{解}{と}ける。\bkfurref{bk:3-10:insu}
\item \R{解}{かい}の\R{公式}{こう|しき}を\R{作}{つく}り、$\sqrt{\phantom{2}}$ の\R{計算}{けい|さん}をして\R{解}{と}ける。\bkfurref{bk:3-10:kousiki}
\item \R{判別式}{はん|べつ|しき}で\R{解}{かい}の\R{個数}{こ|すう}を\R{調}{しら}べ、\R{場面}{ば|めん}に\R{合}{あ}う\R{解}{かい}を\R{選}{えら}べる。\bkfurref{bk:3-10:hanbetsu}
\end{checklist}
\end{furikaeri}

\bkfukidashi{delta}{\R{答}{こた}えが2つ\R{出}{で}たら、どちらも\R{場面}{ば|めん}に\R{合}{あ}うか\R{確}{たし}かめよう。\R{次}{つぎ}の\R{章}{しょう}では、\R{倍々}{ばい|ばい}でふえる\R{量}{りょう}のグラフを\R{見}{み}ていくよ。}
